%\documentclass[11pt,a4paper,twoside]{article}

\usepackage[utf8]{inputenc}

\usepackage[T1,T2A]{fontenc}

\usepackage[english.ancient,russian.ancient]{babel}

\usepackage{graphicx}

\usepackage{array}

\usepackage[doi]{samgtu-bib}

\usepackage{samgtu}

\usepackage{samgtu-name}

\usepackage{mathtools}

\mathtoolsset{showonlyrefs}

\usepackage{desclist}

\usepackage{euscript}

\usepackage{mathrsfs} 

\usepackage{multicol}

\usepackage{mathrsfs}

\usepackage{cite}

\usepackage{cmap}

\usepackage{qrcode}

\allowdisplaybreaks[4]

\usepackage[nodayofweek]{datetime}

\usepackage[thinlines]{easybmat}



\renewcommand{\JournalYear}{2018}

\renewcommand{\JournalVolume}{22}

\renewcommand{\JournalNumber}{4}



\newdate{JournalDateSign}{29}{12}{2018}% Дата подписания Вестника в печать

\renewcommand{\JournalDateSign}{\displaydate{JournalDateSign}}

\newdate{JournalDatePrint}{16}{01}{2019}% Дата выхода Вестника из печати

\renewcommand{\JournalDatePrint}{\displaydate{JournalDatePrint}}

%\renewcommand{\JournalUPL}{16.56} % Усл. печ. л.

%\renewcommand{\JournalUIL}{16.53} % Уч.-изд. л.

\renewcommand{\JournalUPL}{15.24} % Усл. печ. л.

\renewcommand{\JournalUIL}{15.21} % Уч.-изд. л.

\renewcommand{\JournalRegNumber}{220/18} % Регистрационный номер

\renewcommand{\JournalOrderNumber}{2} % Номер заказа



%\renewcommand{\JournalRMonth}{Март} % Месяц выхода Вестника

%\renewcommand{\JournalEMonth}{March} % Месяц выхода Вестника

%\renewcommand{\JournalRMonth}{Июнь} % Месяц выхода Вестника

%\renewcommand{\JournalEMonth}{June} % Месяц выхода Вестника

%\renewcommand{\JournalRMonth}{Сентябрь} % Месяц выхода Вестника

%\renewcommand{\JournalEMonth}{September} % Месяц выхода Вестника

\renewcommand{\JournalRMonth}{Декабрь} % Месяц выхода Вестника

\renewcommand{\JournalEMonth}{December} % Месяц выхода Вестника



\newcommand{\totop}[1]{\raisebox{6pt}[1pt][2pt]{#1}}

\newcommand{\tottop}[1]{\raisebox{12pt}[1pt][2pt]{#1}} 

\newcommand{\totttop}[1]{\raisebox{18pt}[1pt][2pt]{#1}} 

\newcommand{\tottttop}[1]{\raisebox{24pt}[1pt][2pt]{#1}} 

\newcommand{\totttttop}[1]{\raisebox{30pt}[1pt][2pt]{#1}} 

\newcommand{\tobot}[1]{\raisebox{-6pt}[1pt][2pt]{#1}} 

\newcommand{\tobbot}[1]{\raisebox{-12pt}[1pt][2pt]{#1}} 

\newcommand{\tobbbot}[1]{\raisebox{-18pt}[1pt][2pt]{#1}}



\graphicspath{{./00PIC/}}
%
%\usepackage{samgtu-name}

\vsgtu{1600} %registration number

\FirstPage{269}
\LastPage{292}
%
%\pdfcrop
%\draft

\ResearchArticle

%\begin{document}

\hfill \qrcode[hyperlink,height=.8in]{http://doi.org/10.14498/vsgtu1600}
\vspace{-.8in}

\udc{517.956.6}
\msc{35R30,  35M13}

\rutitle{К вопросу о корректности обратных задач для~неоднородного уравнения Гельмгольца}
\rutitleshort{К вопросу о корректности обратных задач\dots}

\entitle{On the question of the correctness of inverse problems for~the~inhomogeneous Helmholtz equation}
\entitleshort{On the question of the correctness of inverse problems \dots}

\ruauthor{К.~Б.~Сабитов$^{1,2}$, Н.~В.~Мартемьянова$^3$}{С\,а\,б\,и\,т\,о\,в~К.~Б., М\,а\,р\,т\,е\,м\,ь\,я\,н\,о\,в\,а~Н.~В.}
\enauthor{K.~B.~Sabitov$^{1,2}$, N.~V.~Martem'yanova$^3$}{S\,a\,b\,i\,t\,o\,v~K.~B., M\,a\,r\,t\,e\,m'\,y\,a\,n\,o\,v\,a~N.~V.}


\ruaffil{\rusamgtu.\and 
\ruaffid{6508}{Самарский государственный социально-педагогический университет,\\
Россия, 443099, Самара, ул. Горького, 65/67}.
\and
\rusgau.
}

\enaffil{\ensamgtu.\and 
\enaffid{6508}{Samara State University of Social Sciences and Education,\\
65/67, M. Gorkiy str., Samara,  443099, Russian Federation}.\and
\ensgau.}



\ruauthorsdetails{2}{% Количество авторов
\fullauthor{\rupid{11101}{Камиль Басирович Сабитов}}
\orcid{0000-0001-9516-2704} % ORCID ID автора 
\regalia{доктор физико-математических  наук, профессор} 
\position{профессор, каф. высшей математики$^1$; профессор, каф. физики, математики и методики обучения$^2$} 
\email{sabitov_fmf@mail.ru}
\fullauthor{\rupid{53362}{Нина Викторовна Мартемьянова}}
\CAuthor % Автор-корреспондент
\orcid{0000-0001-8025-0234} % ORCID ID автора 
\regalia{кандидат физико-математических  наук} 
\position{доцент} 
\dept{каф. высшей математики$^3$} 
\email[b]{ninamartem@yandex.ru}
}


\enauthorsdetails{2}{% Количество авторов
\fullauthor{\enpid{11101}{Kamil B. Sabitov}}
\orcid{0000-0001-9516-2704} % ORCID ID автора 
\regalia{Dr. Phys \& Math. Sci.} 
\position{Professor, Dept. of Higher Mathematics$^1$, Professor, Dept. of Physics, Mathematics and Methods of Teaching$^2$} 
\email{sabitov_fmf@mail.ru}
\fullauthor{\rupid{53362}{Nina V. Martem'yanova}}
\CAuthor % Автор-корреспондент
\orcid{0000-0001-8025-0234} % ORCID ID автора 
\regalia{Cand. Phys \& Math. Sci.}  
\position{Associate Professor} 
\dept{Dept. of Higher Mathematics$^3$} 
\email[b]{ninamartem@yandex.ru}
}


\rukeywords{уравнение Гельмгольца, начально-граничная задача, нелокальные задачи, обратные задачи, единственность, существование, ряд, устойчивость, интегральные уравнения}
\enkeywords{Helmholtz equation, initial-boundary value problem, nonlocal problems, inverse problems, uniqueness, existence, series, stability, integral equations}

\ruabstract{Для уравнения Гельмгольца в прямоугольной области изучены начально"=граничная задача и ее нелокальные модификации, и обратные задачи по его отысканию правой части. Решения прямых задач с нелокальными граничными условиями и обратных задач построены в явном виде как суммы ортогональных рядов по системе собственных функций одномерной спектральной задачи Штурма--Лиувилля. Доказаны соответствующие теоремы единственности решения всех поставленных задач. Установлены достаточные условия на граничные функции, которые гарантируют теоремы существования и устойчивости решения предложенных новых постановок задач.
}

\enabstract{In the rectangular domain, the initial-boundary value problem for the Helmholtz equation and its non-local modifications are studied and the inverse problems for finding its right-hand side are studied. The solutions of direct problems with nonlocal boundary conditions and inverse problems are constructed in explicit form as the sums of orthogonal series in the system of eigenfunctions of the one-dimensional Sturm--Liouville spectral problem. The corresponding uniqueness theorems for the solution of all set problems are proved. Sufficient conditions for boundary functions are established, which are guaranteed by the existence and stability theorems for the solution of the proposed new problem statements.
}

\newdate{sabitova}{10}{01}{2018}
\date{\displaydate{sabitova}}
\newdate{sabitovb}{21}{04}{2018}
\revisiondate{\displaydate{sabitovb}}
\newdate{sabitovc}{11}{06}{2018}
\accepteddate{\displaydate{sabitovc}}

\newdate{sabitove}{27}{06}{2018}
\renewcommand{\JournalDataOnline}{\displaydate{sabitove}}


%\ForPeerReview
\makerutitle




\Section[n]{Введение}
В связи с изучением обратной задачи по  определению правой части уравнения смешанного эллиптико"=гиперболического типа, например, известного уравнения Лаврентьева"--~Бицадзе \hypertarget{(ast)}{}
$$
u_{xx}+(\mathop{\rm sgn} y)u_{yy}=\Phi(x,y)
%\label{(ast)}
\eqno{(*)}
$$
в прямоугольной области $G=\{(x,y) \mid  0< x<l,\,-\alpha<y<\beta\}$, где
$$
\Phi(x,y)=
\begin{cases}
      \Phi_1(x,y)=f_1(x)g_1(y), & y>0, \\
      \Phi_2(x,y)=f_2(x)g_2(y), & y<0,
    \end{cases}
$$
возникает необходимость исследования обратных задач по отысканию правых частей уравнений 
$u_{xx}+u_{yy}=\Phi_1(x,y)$ 
в~$G_+=G\cap\{y>0\}$ и $u_{xx}-u_{yy} \hm=\Phi_2(x,y)$ в~$G_-=G\cap\{y<0\}$. 
Ранее в наших работах \cite{sab:1, sab:2, sab:3, sab:4, sab:5, sab:6} были изучены обратные задачи по определению функций $u(x,y)$ и $f_1(x)=f_2(x)=f(x)$ и $u(x,y)$, $f_1(x)$ и $f_2(x)$, $f_1(x)\neq f_2(x)$ для уравнения  (\hyperlink{(ast)}{$*$}) в~области $G$ при $g_1(y)=g_2(y)\equiv1$ с~локальными и~нелокальными граничными условиями. 
Эти исследования не позволяют пока построить теорию обратных задача для уравнения (\hyperlink{(ast)}{$*$}). 
В~связи с~этим важно предварительно исследовать обратные задачи для уравнения (\hyperlink{(ast)}{$*$}) в~областях эллиптичности и~гиперболичности.

Отметим, что задачи оптимизации и связанные с ними обратные задачи, например, задача оптимизации сопла Лаваля, являются предметом исследований работ \cite{sab:7, sab:8}.

Обратные задачи разного плана для классических уравнений математической физики и для более общих дифференциальных уравнений в частных производных различных типов изучены достаточно полно; усилиями многих математиков создана теория обратных задач (см. \cite{sab:9, sab:10, sab:11, sab:12, sab:13, sab:14, sab:15, sab:16} и приведенную там обширную библиографию).

В~данной работе мы исследуем прямые и~обратные задачи для уравнений Лапласа и Гельмгольца.
К нашим исследованиям близки работы \cite{sab:17, sab:18, sab:19, sab:20}, где для уравнения Пуассона и~более общего уравнения эллиптического типа изучены обратные задачи по отысканию функций $u(x,y)$ и $f_1(x)$ в~прямоугольной области и~областях более общего вида с~заданием граничного условия на всей границе области и дополнительной информации о~функции $u(x,y)$ внутри области и доказаны теоремы единственности и существования в классах гладких функций. 

Отметим также работы \cite{sab:21, sab:22, sab:23}, в~которых изучена задача по определению сомножителя правой части операторного уравнения эллиптического типа, где установлен критерий единственности и~найдено необходимое и~достаточное условие разрешимости поставленной задачи.

В отличие от этих работ, в~данной статье решения поставленных обратных задач построены в~явном виде как суммы рядов, доказаны соответствующие теоремы единственности, существования и~устойчивости.

В связи с этим п.~\hyperlink{sab:p1}{1} данной работы посвящен изучению начально"=граничной задачи, имеющей важное значение в геофизике \cite[с.~131]{sab:13}, \cite[с.~254]{sab:16}, и~двух ее нелокальных модификаций для уравнения Гельмгольца \hypertarget{(astast)}{}
$$
u_{xx}+u_{yy}+cu=F(x,y)=f(x)g(y)
%\label{(astast)}
\eqno{(**)}
$$
в прямоугольной области $G_+$, которые являются корректными по Адамару и~будут использованы в последующем при постановке обратных задач.

В п.~\hyperlink{sab:p2}{2} изучены обратные задачи по отысканию правой части уравнения~(\hyperlink{(astast)}{$**$}).

В пп.~\hyperlink{sab:p2.1}{2.1} рассмотрен случай, когда $g(y)\equiv1$ и неизвестными являются функции $u(x,y)$ и~$f(x)$. 
Решение построено в~виде сумм ортогональных рядов и доказаны теоремы единственности, существования и~устойчивости решения.

В пп.~\hyperlink{sab:p2.2}{2.2} изучен случай, когда $g(y)\not\equiv1$ и неизвестными также являются функции $u(x,y)$ и~$f(x)$. 
Установлен критерий единственности решения поставленной обратной задачи. 
Аналогично пп.~\hyperlink{sab:p2.1}{2.1}, решение построено в виде сумм ортогональных рядов и~доказаны теоремы существования и устойчивости решения поставленной обратной задачи.

В пп.~\hyperlink{sab:p2.3}{2.3} изучается обратная задача, в которой неизвестными являются функции $u(x,y)$ и~$g(y)$. 
На основании результатов п.~\hyperlink{sab:p1}{1} относительно функции $g(y)$ получено интегральное уравнение Фредгольма второго рода, из которого при определенных условиях на данные следует однозначная разрешимость обратной задачи.


\smallskip

\hypertarget{sab:p1}{}
\Section{Начально-граничная задача и ее нелокальные модификации}
Рассмотрим неоднородное уравнение Гельмгольца
\begin{equation}\label{o1}
Lu=u_{xx}+u_{yy}+cu=F(x,y),
\end{equation}
где $c={\rm const}<(\pi/l)^2$, в прямоугольной области $D=\{(x,y)\mid 0<x<l,$ ${0<y<\beta}\}$ и следующую начально"=граничную задачу.

\smallskip

\hypertarget{sab:task1}{}

{\small\sc Задача 1.} Найти функцию $u(x,y)$, удовлетворяющую условиям
\begin{gather}
\label{o2}
u(x,y)\in C^{2}(D)\cap C^{1}(D\cup\{0\leq x\leq l, \,y=0\})\cap C(\overline{D}),
\\
\label{o3}
Lu(x,y)\equiv F(x,y), \quad  (x,y)\in D,
\\ 
 \label{o4}
u(0,y)=u(l,y)=0, \quad  0\leq y\leq \beta,
\\
\label{o5}
u(x,0)=\varphi(x), \quad  0\leq x\leq l,
\\
\label{o6}
u_y(x,0)=\psi(x), \quad  0\leq x\leq l,
\end{gather}
где $F(x,y)$, $\varphi(x)$, $\psi(x)$ "--- заданные достаточно гладкие функции.

\smallskip

Отметим, что к этой задаче сводятся двумерные задачи Коши для уравнения Лапласа, связанные с изучением стационарных электростатических или магнитных полей с целью поиска полезных ископаемых \cite[с.~131]{sab:13}, \cite[с.~254]{sab:16}.

Рассматриваемая задача \eqref{o2}--\eqref{o6} некорректна по Адамару. 
Пусть ${F(x,y){\equiv}0}$. 
Тогда решением задачи \hyperlink{sab:task1}{1} является функция
\begin{equation}\label{o7}
u_k(x,y)=\frac{1}{k}\sin\mu_kx \Bigl(a\ch\lambda_ky+\frac{b}{\lambda_k}\sh\lambda_ky\Bigr)
\end{equation}
с начальными функциями
\begin{equation}\label{o8}
\varphi_k(x)=\frac{a}{k}\sin\mu_kx,\quad  
\psi_k(x)=\frac{b}{k}\sin\mu_kx,
\end{equation}
$a$, $b$ "--- вещественные постоянные, $k\in\mathbb{N}$.

Из \eqref{o7} и \eqref{o8} следует, что с ростом $k$ начальные функции $\varphi_k(x)$ и $\psi_k(x)$ стремятся к~нулю, а $u_k(x,y)$ неограниченно растет, т.~е. решение задачи \hyperlink{sab:task1}{1} неустойчиво.

\smallskip

\hypertarget{sab:th1}{}

{\small\sc Теорема 1.} \textit{Если существует решение задачи {\em \eqref{o2}--\eqref{o6},} удовлетворяющее условиям
\begin{equation}
\label{o9}
\lim_{x\rightarrow0+0}u_x\sin x=\lim_{x\rightarrow l-0}u_x\sin\frac{\pi x}{l}=0,
\end{equation}
то оно единственно}.

\smallskip

\textit{Д\,о\,к\,а\,з\,а\,т\,е\,л\,ь\,с\,т\,в\,о.} 
Пусть существует решение задачи \eqref{o2}--\eqref{o6}, удовлетворяющее условиям \eqref{o9}. 
Следуя \cite{sab:3}, введем  функции
\begin{equation}\label{o10}
u_k(y)=\frac{2}{l}\int _{0}^{l} u(x,y)X_k(x)\,dx, \quad k\in\mathbb{N},
\end{equation}
где $X_k(x)=\sin\mu_kx$, $\mu_k=\pi k/l$ "--- собственные функции спектральной задачи
$$
X''(x)+\mu^2X(x)=0,\quad  X(0)=X(l)=0,
$$
которые образуют ортогональную и~полную систему в~пространстве $L_2[0,l]$.

Рассмотрим
\begin{equation}\label{o11}
u_{k,\varepsilon}(y)=\frac{2}{l}\int _{\varepsilon}^{l-\varepsilon} u(x,y)X_k(x)\,dx,
\end{equation}
где $\varepsilon>0$ "--- достаточно малое число. 
Дифференцируя равенство \eqref{o11} дважды и используя уравнение \eqref{o1}, имеем
\begin{multline}
u''_{k,\varepsilon}(y)=\frac{2}{l}\int _{\varepsilon}^{l-\varepsilon} u_{yy}(x,y)X_k(x)\,dx= \frac{2}{l}\int _{\varepsilon}^{l-\varepsilon} \bigl[F(x,y)-u_{xx}-cu\bigr]X_k(x)\,dx=
\\
\label{o12}
=\frac{2}{l}\int _{\varepsilon}^{l-\varepsilon} F(x,y)X_k(x)\,dx-cu_{k,\varepsilon}-\frac{2}{l}\int _{\varepsilon}^{l-\varepsilon} u_{xx}(x,y)X_k(x)\,dx.
\end{multline}
В последнем интеграле правой части \eqref{o12} после интегрирования по частям два раза и перехода к пределу при $\varepsilon\rightarrow0$ с~учетом условий \eqref{o4}--\eqref{o9} получим, что функции \eqref{o10} удовлетворяют дифференциальному уравнению
\begin{equation}\label{o13}
u''_k(y)-\lambda^2_ku_k(y)=F_k(y).
\end{equation}
Здесь $\lambda^2_k=\mu^2_k-c>0$,
\begin{equation}\label{o14}
F_k(y)=\frac{2}{l}\int _{0}^{l} F(x,y)X_k(x)\,dx.
\end{equation}
Общее решение уравнения \eqref{o13} определим в виде
\begin{equation}\label{o15}
u_k(y)=a_ke^{\lambda_ky}+b_ke^{-\lambda_ky}-\frac{1}{2\lambda_k}F_{1k}(y)-\frac{1}{2\lambda_k}F_{2k}(y),
\end{equation}
удобном для дальнейших исследований. Здесь $a_k$ и $b_k$ "--- произвольные постоянные,
$$
F_{1k}(y)=\int _{y}^{\beta} F_k(t)e^{-\lambda_k(t-y)}\,dt,\quad  
F_{2k}(y)=\int _{0}^{y} F_k(t)e^{-\lambda_k(y-t)}\,dt.
$$

Для нахождения $a_k$ и $b_k$ воспользуемся граничными условиями \eqref{o5} и \eqref{o6}:
\begin{equation}\label{o16}
u_k(0)=\frac{2}{l}\int _{0}^{l} u(x,0)\sin\mu_kx\,dx=\frac{2}{l}\int _{0}^{l} \varphi(x)\sin\mu_kx\,dx=\varphi_k,
\end{equation}
\begin{equation}\label{o17}
u'_k(0)=\frac{2}{l}\int _{0}^{l} u_y(x,0)\sin\mu_kx\,dx=\frac{2}{l}\int _{0}^{l} \psi(x)\sin\mu_kx\,dx=\psi_k.
\end{equation}
Теперь, удовлетворяя функции \eqref{o15} условиям \eqref{o16} и \eqref{o17}, получим систему
\begin{equation}\label{o18}
\begin{cases}
      \displaystyle a_k+b_k=\varphi_k+\frac{1}{2\lambda_k}F_{1k}(0), \\[2mm]
      \displaystyle a_k-b_k=\frac{\psi_k}{\lambda_k}+\frac{1}{2\lambda_k}F_{1k}(0),
    \end{cases}
\end{equation}
из которой определяются $a_k$ и $b_k$:
$$
a_k=\frac{\varphi_k+\psi_k/\lambda_k}{2}+\frac{1}{2\lambda_k}F_{1k}(0),\quad b_k=\frac{\varphi_k-\psi_k/\lambda_k}{2}.
$$
Подставляя найденные значения $a_k$ и $b_k$ в \eqref{o15}, получим явный вид функций~$u_k(y)$:
\begin{equation}\label{o19}
u_k(y)=\varphi_k\ch\lambda_ky+\frac{\psi_k}{\lambda_k}\sh\lambda_ky+ \frac{1}{2\lambda_k}e^{\lambda_ky}F_{1k}(0)
-\frac{1}{2\lambda_k}F_{1k}(y)- \frac{1}{2\lambda_k}F_{2k}(y).
\end{equation}

Пусть теперь $\varphi(x)\equiv0$, $\psi(x)\equiv0$, $F(x,y)\equiv0$. 
Тогда все $\varphi_k=\psi_k\equiv0$, $F_k(t)\equiv0$ и из \eqref{o19} и \eqref{o10} следует, что при всех $k\in\mathbb{N}$
\begin{equation}\label{o20}
\int _{0}^{l}u(x,y)\sin\mu_kx\,dx=0.
\end{equation}
Из \eqref{o20} в силу полноты системы $\bigl\{\sin\mu_kx\bigr\}_{k\geq1}$ в $L_2[0,l]$ следует, 
что ${u(x,y)=0}$ почти всюду на $[0,l]$ при любом $y\in[0,\beta]$. 
Поскольку в силу условия \eqref{o2} функция $u(x,y)$ непрерывна на $\overline{D}$, она тождественно равна нулю  на~$\overline{D}$.  $\square$

\bigskip


{\small\sc  Замечание 1.} Отметим, что в работе \cite[с. 254--256]{sab:16}  приведено доказательство теоремы единственности решения задачи~\hyperlink{sab:task1}{1} в~классе функций из $C^4(\overline{D})$, представимых в~виде ряда
$$
u(x,y)=\sum\limits_{k=1}^{\infty} u_k(y)\sin{\pi kx},
$$ 
где  $u_k(y)\in C^4[0,1]$, $l=\beta=1$.
В~теореме~\hyperlink{sab:th1}{1} условия \eqref{o9} означают, что производная $u_x$ вблизи отрезков $x=0$ и $x=l$ может иметь особенности порядка меньше единицы.

\smallskip


Формально решение задачи~\hyperlink{sab:task1}{1} можно построить в виде суммы ряда
\begin{equation}\label{o21}
u(x,y)=\sum\limits_{k=1}^{\infty} u_k(y)\sin\mu_kx,
\end{equation}
где коэффициенты $u_k(y)$ однозначно определены по формуле \eqref{o19}.

Поскольку система  $\bigl\{\sin\mu_kx \bigr\}_{k\geq1}$ образует ортогональный базис в $L_2[0,l]$, ряд \eqref{o21} сходится в $L_2[0,l]$ при любом фиксированном $y\in[0,\beta]$.
Доказать равномерную сходимость этого ряда на $\overline{D}$ не удастся, так как  коэффициенты $u_k(y)$ при больших $k$ растут по закону экспоненты $e^{\lambda_ky}$. 
Поэтому доказательство существования решения задачи~\hyperlink{sab:task1}{1} в указанном классе функций \eqref{o2} и \eqref{o9} остается под большим вопросом.

Возникает вопрос: можно ли, не меняя структуру решения, поправить постановку задачи \eqref{o2}--\eqref{o6} так, чтобы измененная задача стала корректной?

Здесь рассмотрим два варианта. В первом варианте предлагается оставить граничные условия \eqref{o4} и~\eqref{o5} без изменения, а условие \eqref{o6} заменить нелокальным граничным условием
\begin{equation}\label{o29}
u_y(x,0)-u_y(x,\beta)=\widetilde{\psi}(x),\quad  0\leq x\leq l.
\end{equation} 

\hypertarget{sab:task2}{}

\noindent
Эту задачу назовем {{\small \sc Задачей 2}}.
В этом случае решение задачи~\hyperlink{sab:task2}{2} ищется в~классе функций \hypertarget{sab:(2')}{}
$$
u(x,y)\in C^{2}(D)\cap C^{1}(D\cup\{y=0,\, 0\leq x\leq l\}\cup\{y=\beta,\,0\leq x\leq l\})\cap C(\overline{D}). \eqno{(2')}
$$

Во втором варианте предлагается в задаче~\hyperlink{sab:task1}{1} заменить условие \eqref{o5} на нелокальное условие
\begin{equation}
\label{o30}
u(x,0)-u(x,\beta)=\widetilde{\varphi}(x),\quad  0\leq x\leq l.
\end{equation}

\hypertarget{sab:task2}{}

\noindent
Эту задачу назовем {\small \sc Задачей 3}. 
В~условиях \eqref{o29} и~\eqref{o30} функции $\widetilde{\psi}(x)$ и $\widetilde{\varphi}(x)$ считаются известными и~достаточно гладкими.

Решение задачи (\hyperlink{sab:(2')}{$2'$}), \eqref{o3}--\eqref{o5} и \eqref{o29} (задачи~\hyperlink{sab:task2}{2}), аналогично решению задачи~\hyperlink{sab:task1}{1}, можно построить в виде суммы ряда
\begin{equation}\label{o31}
u(x,y)=\sum\limits_{k=1}^{\infty} u_k(y)\sin\mu_kx.
\end{equation}

Действительно, вводя функции \eqref{o10} относительно $u_k(y),$ получаем уравнение \eqref{o13}, общее решение которого определяется по формуле \eqref{o15}. 
Для определения неизвестных коэффициентов $a_k$ и $b_k$ воспользуемся граничными условиями \eqref{o5} и \eqref{o29}, т.~е. условиями \eqref{o16} и
\begin{multline}\label{o310}
u'_k(0)-u'_k(\beta)=\frac{2}{l}\int _{0}^{l}\bigl[u_y(x,0)-u_y(x,\beta)\bigr]\sin\mu_kx\,dx=
\\
=
 \frac{2}{l}\int _{0}^{l}\widetilde{\psi}(x)\sin\mu_kx\,dx=\widetilde{\psi}_k.
\end{multline}
Удовлетворив функции \eqref{o15} граничным условиям \eqref{o16} и \eqref{o310}, найдем искомые коэффициенты:
\begin{gather}
a_k=\frac{1}{2(\ch\lambda_k\beta-1)}
\Bigl[\varphi_k\bigl(e^{-\lambda_k\beta}-1\bigr)-\frac{\widetilde{\psi}_k}{\lambda_k}+ \frac{F_{1k}(0)}{2\lambda_k}
\bigl(e^{-\lambda_k\beta}-2\bigr)- \frac{F_{2k}(\beta)}{2\lambda_k}
\Bigr],
\\
b_k=\frac{1}{2(\ch\lambda_k\beta-1)}
\Bigl[\frac{\widetilde{\psi}_k}{\lambda_k}-\varphi_k\bigl(1-e^{\lambda_k\beta}\bigr)+ \frac{e^{\lambda_k\beta}F_{1k}(0)}{2\lambda_k}+\frac{F_{2k}(\beta)}{2\lambda_k}
\Bigr].
\end{gather}
Подставляя найденные значения $a_k$ и $b_k$ в формулу \eqref{o15}, получим явный вид искомых функций:
\begin{multline}
u_k(y)=\varphi_k\frac{\ch\lambda_k(\beta-y)-\ch\lambda_ky}{(\ch\lambda_k\beta-1)}- \widetilde{\psi}_k\frac{\sh\lambda_ky}{\lambda_k(\ch\lambda_k\beta-1)}+
\\
\label{o32}
+\frac{F_{1k}(0)[\ch\lambda_k(\beta-y)-\ch\lambda_ky-\sh\lambda_ky]}{2\lambda_k(\ch\lambda_k\beta-1)}- \\
- \frac{1}{2\lambda_k}F_{1k}(y)-\frac{F_{2k}(\beta)\sh\lambda_ky}{2\lambda_k(\ch\lambda_k\beta-1)}- \frac{1}{2\lambda_k}F_{2k}(y).
\end{multline}

\smallskip

\hypertarget{sab:lemma1}{}

{\small \sc Лемма 1.} \textit{При любом $y\in[0,\beta]$ для функций \eqref{o32} справедливы оценки
\begin{gather}\label{o33}
|u_k(y)|\leq C_1\Bigl(|\varphi_k|+\frac{1}{k}|\widetilde{\psi}_k|+\frac{1}{k^{3/2}}\|F_k(y)\|_{L_2}\Bigr),
\\
\label{o35}
|u'_k(y)|\leq C_2\Bigl(k|\varphi_k|+|\widetilde{\psi}_k|+\frac{1}{k^{1/2}}\|F_k(y)\|_{L_2}\Bigr),
\\
\label{o35}
|u''_k(y)|\leq C_3\Bigl(k^2|\varphi_k|+k|\widetilde{\psi}_k|+k^{1/2}\|F_k(y)\|_{L_2}\Bigr).
\end{gather}
Здесь и ниже $C_i$ "---  положительные постоянные$,$ зависящие$,$ вообще говоря$,$ от данных задачи$,$ $\beta,$ $c$ и $l;$ $\|F_k(y)\|_{L_2}=\|F_k(y)\|_{L_2[0,\beta]}.$}

\smallskip

\textit{Д\,о\,к\,а\,з\,а\,т\,е\,л\,ь\,с\,т\,в\,о}. Отметим, что
\begin{gather}
\frac{|\ch\lambda_k(\beta-y)-\ch\lambda_ky|}{\ch\lambda_k\beta-1}\leq1,\quad
0\leq\frac{\sh\lambda_ky}{\ch\lambda_k\beta-1}\leq\frac{1+e^{-\lambda_1\beta}}{1-e^{-\lambda_1\beta}},
\\
0<\frac{\sh\lambda_k(\beta-y)+\sh\lambda_ky}{\ch\lambda_k\beta-1}\leq\frac{1+e^{-\lambda_1\beta}}{1-e^{-\lambda_1\beta}},\quad 
0<\frac{\ch\lambda_ky}{\ch\lambda_k\beta-1}\leq\frac{1+e^{-2\lambda_1\beta}}{(1-e^{-\lambda_1\beta})^2},
\\
|F_{1k}(y)|\leq\frac{\|F_k(y)\|_{L_2}}{\sqrt{2\lambda_k}},\quad  
|F_{2k}(y)|\leq\frac{\|F_k(y)\|_{L_2}}{\sqrt{2\lambda_k}}.
\end{gather}


Тогда оценки \eqref{o33}--\eqref{o35} непосредственно следуют из формулы \eqref{o32}.~$\square$

\smallskip

На основании оценок \eqref{o33}--\eqref{o35} нетрудно доказать следующие утверждения.

\smallskip

\hypertarget{sab:th2}{}

{\small \sc Теорема 2.} \textit{Если $\varphi(x)\in C^3[0,l],$ 
$\varphi(0)=\varphi''(0)=0,$ 
$\widetilde{\psi}(x)\in C^2[0,l],$ 
$\widetilde{\psi}(0)\hm =\widetilde{\psi}(l)=0;$ 
$F(x,y),$ $F'_x(x,y),$ $F''_{xx}(x,y)\in C(\overline{D}),$ 
$F(0,y)=F(l,y)=0$ при $0\leq y\leq \beta,$ 
то существует единственное решение $u(x,y)$ задачи \em (\hyperlink{sab:(2')}{$2'$}),  \eqref{o3}--\eqref{o5} \em и \em \eqref{o29}, \em которое определяется как сумма ряда \em \eqref{o31}, \em а~коэффициенты "--- по формуле \em \eqref{o32}, \em при этом $u(x,y)\in C^2(\overline{D}).$}

\smallskip

\textit{Д\,о\,к\,а\,з\,а\,т\,е\,л\,ь\,с\,т\,в\,о}. В силу леммы~\hyperlink{sab:lemma1}{1} ряд \eqref{o31} и ряды, полученные из него почленным дифференцированием до второго порядка включительно, в~замкнутой области $\overline{D}$ мажорируются числовым рядом
$$
C_4\sum_{k=1}^{+\infty} k^2|\varphi_k|+k|\widetilde{\psi}_k|+k^{1/2}\|F_k(y)\|_{L_2},
$$
который в силу условий теоремы сходится \cite[с.~90--94]{sab:24}. 
Тогда в силу признака Вейерштрасса указанные ряды сходятся равномерно на $\overline{D}$. 
Поэтому сумма ряда \eqref{o31} принадлежит пространству $C^2(\overline{D})$ и~удовлетворяет граничным условиям \eqref{o4}, \eqref{o5} и \eqref{o29}. $\square$

\smallskip

\hypertarget{sab:th3}{}

{\small \sc Теорема 3.} \textit{Для решения задачи \em (\hyperlink{sab:(2')}{$2'$}), \eqref{o3}--\eqref{o5} \em и \em  \eqref{o29} \em справедливы оценки  устойчивости}\/:
\begin{gather}\label{o36}
\|u(x,y)\|_{L_2[0,l]}\leq C_5 
\bigl(\|\varphi(x)\|_{L_2[0,l]}+\|\widetilde{\psi}(x)\|_{L_2[0,l]}+\|F(x,y)\|_{L_2(D)}\bigr),
\\
\label{o37}
\|u(x,y)\|_{C(\overline{D})}\leq C_6\bigl(\|\varphi'(x)\|_{C[0,l]}+\|\widetilde{\psi}(x)\|_{C[0,l]}+\|F(x,y)\|_{C(\overline{D})}\bigr).
\end{gather}

\smallskip 

 
\textit{Д\,о\,к\,а\,з\,а\,т\,е\,л\,ь\,с\,т\,в\,о.} В силу оценки \eqref{o33} на основании формулы \eqref{o31} имеем
\begin{multline}
\|u(x,y)\|^2_{L_2[0,l]}=\frac{l}{2}\sum_{k=1}^{\infty} u^2_k(y)\leq 
\frac{3lC^2_1}{2}\sum_{k=1}^{\infty} 
\Bigl(|\varphi_k|^2+|\widetilde{\psi}_k|^2+\frac{1}{k^3}\|F_k(y)\|^2_{L_2}\Bigr)\leq
\\
\leq C^2_5\bigl(\|\varphi\|^2_{L_2[0,l]}+\|\widetilde{\psi}\|^2_{L_2[0,l]}+\|F(x,y)\|^2_{L_2(D)}\bigr).
\end{multline}
Отсюда следует справедливость оценки \eqref{o36}.

Пусть $(x,y)$ "--- произвольная точка области $\overline{D}$. 
Тогда из соотношений \eqref{o31} и \eqref{o33} будем иметь
\begin{equation}
\label{o38}
|u(x,y)|\leq\sum_{k=1}^{\infty}|u_k(y)|\leq C_1\sum_{k=1}^{\infty} 
\Bigl(|\varphi_k|+\frac{1}{k}|\widetilde{\psi}_k|+\frac{1}{k^{3/2}}\|F_k(y)\|_{L_2}\Bigr).
\end{equation}
Теперь, используя представление $\varphi_k=\varphi_k^{(1)}/\mu_k$, где 
$$\varphi_k^{(1)}=\cfrac{2}{l}\int _{0}^{l}\varphi'(x)\cos\mu_kx\,dx,$$ оценку
$$
\|F_k(y)\|_{L_2}=\|F_k(y)\|_{L_2[0,\beta]}\leq\sqrt{\frac{2}{l}}\,\|F(x,y)\|_{L_2(D)}\leq \sqrt{2\beta}\max_{\overline{D}}|F(x,y)|
$$
и неравенство Коши"--~Буняковского из \eqref{o38}, получаем
\begin{multline}
|u(x,y)|\leq\widetilde{C}_1\sum\limits_{k=1}^{\infty}
\Bigl(\frac{1}{k}|\varphi^{(1)}_k|+ \frac{1}{k}|\widetilde{\psi}_k|+\frac{\|F(x,y)\|_{C(\overline{D})}}{k^{3/2}}\Bigr)\leq
\\
\leq\widetilde{C}_1\sum\limits_{k=1}^{\infty}
\Bigl(\frac{1}{k^2}\Bigr)^{1/2} 
\biggl[
\biggl(\sum\limits_{k=1}^{\infty}|\varphi^{(1)}_k|^2\biggr)^{1/2}+ 
\biggl(\sum\limits_{k=1}^{\infty}|\widetilde{\psi}_k|^2\biggr)^{1/2}\biggr]+ \hspace{2cm}
\\
\hspace{5cm}+
\widetilde{C}_1\|F(x,y)\|_{C(\overline{D})}\sum\limits_{k=1}^{\infty}\frac{1}{k^{3/2}}=
\\
=\widetilde{C}_1\frac{\pi}{\sqrt{6}} 
\bigl(\|\varphi'(x)\|_{L_2[0,l]}+\|\widetilde{\psi}\|_{L_2[0,l]}\bigr)+ \widetilde{C}_1\widetilde{C}_2\|F(x,y)\|_{C(\overline{D})}\leq
\\
\label{o39}
\leq\widetilde{C}_1\frac{\pi\sqrt{l}}{\sqrt{6}}\left(\|\varphi'(x)\|_{C[0,l]}+\|\widetilde{\psi}\|_{C[0,l]}\right)+ \widetilde{C}_1\widetilde{C}_2\|F(x,y)\|_{C(\overline{D})},
\end{multline}
где 
$$\widetilde{C}_1=C_1\max\{1,l/\sqrt{\pi}, \sqrt{2\beta}\},\quad  
\sum\limits_{k=1}^{\infty} \frac 1{k^2}= \frac{\pi^2}6,\quad  \sum\limits_{k=1}^{\infty} \frac1 {k^{3/2}}=\widetilde{C}_2>0.$$
Из неравенства \eqref{o39} непосредственно следует оценка \eqref{o37}.~$\square$

\smallskip


Аналогично задачам \hyperlink{sab:task1}{1} и \hyperlink{sab:task2}{2}, построим решение задачи \eqref{o2}--\eqref{o4}, \eqref{o6} и \eqref{o30} (задачи\hypertarget{sab:task3}{~}3) как сумму ряда
\begin{equation}
\label{o40}
u(x,y)=\sum\limits_{k=1}^{\infty}u_k(y)\sin\mu_kx.
\end{equation}
Здесь $u_k(y)$ определяется по формуле
\begin{multline}
u_k(y)=-\widetilde{\varphi}_k\frac{\ch\lambda_ky}{\ch\lambda_k\beta-1}- \frac{\psi_k}{\lambda_k}\frac{\sh\lambda_k(\beta-y)+\sh\lambda_ky}{\ch\lambda_k\beta-1}-
\\
\label{o41}
- \frac{F_{1k}(0)}{2\lambda_k}
\frac{\sh\lambda_k(\beta-y)+\sh\lambda_ky+\ch\lambda_ky }{\ch\lambda_k\beta-1}- \frac{1}{2\lambda_k}F_{1k}(y)+
\\ 
+\frac{F_{2k}(\beta)}{2\lambda_k}\frac{\ch\lambda_ky}{\ch\lambda_k\beta-1}- \frac{1}{2\lambda_k}F_{2k}(y),
\end{multline}
где
$$
\widetilde{\varphi}_k=\cfrac{2}{l}\int _{0}^{l}\widetilde{\varphi}(x)\sin\mu_kx\,dx.
$$

\smallskip

\hypertarget{sab:lemma2}{}

{\small \sc Лемма 2.} \textit{При любом $y\in[0,\beta]$ справедливы оценки}
\begin{gather}
|u_k(y)|\leq C_7\Bigl(|\widetilde{\varphi}_k|+\frac{|\psi_k|}{k}+\frac{1}{k^{3/2}}\|F_k(y)\|_{L_2[0,\beta]}\Bigr),
\\
|u'_k(y)|\leq C_8\Bigl(k|\widetilde{\varphi}_k|+|\psi_k|+\frac{1}{k^{1/2}}\|F_k(y)\|_{L_2[0,\beta]}\Bigr),
\\
|u''_k(y)|\leq C_9\bigl(k^2|\widetilde{\varphi}_k|+k|\psi_k|+k^{1/2}\|F_k(y)\|_{L_2[0,\beta]}\bigr).
\end{gather}

\smallskip

\textit{Д\,о\,к\,а\,з\,а\,т\,е\,л\,ь\,с\,т\,в\,о\/} аналогично обоснованию леммы \hyperlink{sab:lemma1}{1}.

\smallskip

\hypertarget{sab:th4}{}

{\small \sc Теорема 4.} \textit{Если функция $F(x,y)$ удовлетворяет условиям теоремы {\em \hyperlink{sab:th2}{2}}, $\widetilde{\varphi}(x)\,{\in}\, C^3[0,l],$
$\widetilde{\varphi}(0)=\widetilde{\varphi}(l)=\widetilde{\varphi}''(0)=\widetilde{\varphi}''(l)=0,$ 
$\psi(x)\in C^2[0,l],$ ${\psi(0)=\psi(l)=0},$ то существует единственное решение $u(x,y)$ задачи {\em \eqref{o2}--\eqref{o4}, \eqref{o6}} и {\em \eqref{o30},} и это решение определяется рядом {\em \eqref{o40},} у которого коэффициенты находятся по формуле {\em \eqref{o41},} при этом $u(x,y)\in C^2(\overline{D}).$}

\smallskip

\textit{Д\,о\,к\,а\,з\,а\,т\,е\,л\,ь\,с\,т\,в\,о\/}  следует из установленных в лемме \hyperlink{sab:lemma2}{2} оценок.

\smallskip

Для этой задачи имеют место оценки  устойчивости, аналогичные оценкам \eqref{o36} и~\eqref{o37}.




\bigskip
\textbf{2. Обратные задачи по отысканию правой части.}
Здесь для уравнения \eqref{o1} исследуем обратные задачи по определению правой части, т.е. функции $F(x,y)=f(x)g(y)$.

\smallskip

\hypertarget{sab:p2.1}{}

\textbf{2.1.} В начале рассмотрим случай, когда $g(y)\equiv1$.

\smallskip

\hypertarget{sab:task4}{}

{\small\sc Задача 4.} Найти функции $u(x,y)$ и $f(x)$, удовлетворяющие условиям \eqref{o2}--\eqref{o6} и, кроме того,
\begin{gather}
\label{o42}
u(x,\beta)=h(x),\quad  0\leq x\leq l,
\\
\label{o43}
f(x)\in C(0,l)\cap L_2(0,l),
\end{gather}
где $\varphi(x)$, $\psi(x)$ и $h(x)$ "--- заданные достаточно гладкие функции, условие \eqref{o6} здесь является дополнительным для определения функции $f(x)$.

\smallskip

\hypertarget{sab:th5}{}

{\small \sc Теорема 5.} \textit{Если существует решение задачи {\em \hyperlink{sab:task4}{4},} удовлетворяющее условиям {\em \eqref{o9},} то оно единственно}.

\smallskip


\textit{Д\,о\,к\,а\,з\,а\,т\,е\,л\,ь\,с\,т\,в\,о\/.}  Аналогично доказательству теоремы~\hyperlink{sab:th1}{1}, введем функции \eqref{o10}, которые удовлетворяют уравнению
\begin{equation}
\label{o44}
u''_k(y)-\lambda^2_ku_k(y)=f_k,
\end{equation}
где
\begin{equation}
\label{o45}
f_k=\frac{2}{l}\int _{0}^{l}f(x)\sin\mu_kx\,dx.
\end{equation}
Общее решение уравнения \eqref{o44} определяется по формуле
\begin{equation}
\label{o46}
u_k(y)=a_ke^{\lambda_ky}+b_ke^{-\lambda_ky}-\frac{f_k}{\lambda^2_k},
\end{equation}
где $a_k$, $b_k$ и $f_k$ "--- неизвестные постоянные. 
Для их нахождения воспользуемся граничными условиями \eqref{o16}, \eqref{o17} и
\begin{equation}\label{o47}
u_k(\beta)=\frac{2}{l}\int _{0}^{l}u(x,\beta)\sin\mu_kx\,dx= \frac{2}{l}\int _{0}^{l}h(x)\sin\mu_kx\,dx=h_k.
\end{equation}
Удовлетворяя функции \eqref{o46} граничным условиям \eqref{o16}, \eqref{o17} и \eqref{o47}, получим систему
\begin{equation}
\label{o48}
\begin{cases}
      \displaystyle a_k+b_k-\frac{f_k}{\lambda^2_k} =\varphi_k, \\
      \displaystyle a_k-b_k =\frac{\psi_k}{\lambda_k},\\
      \displaystyle a_ke^{\lambda_k\beta}+b_ke^{-\lambda_k\beta}-\frac{f_k}{\lambda^2_k} =h_k,
    \end{cases}
\end{equation}
единственное решение которой определяется формулами
\begin{gather}
\label{o49}
a_k=\frac{h_k-\varphi_k}{2(\ch\lambda_k\beta-1)}- \frac{\psi_k(1-e^{-\lambda_k\beta})}{2\lambda_k(\ch\lambda_k\beta-1)},
\\
\label{o51}
b_k=\frac{\psi_k(1-e^{\lambda_k\beta})}{2\lambda_k(\ch\lambda_k\beta-1)}- \frac{\varphi_k-h_k}{2(\ch\lambda_k\beta-1)},
\\
\label{o51}
\frac{f_k}{\lambda_k^2}=\frac{h_k}{\ch\lambda_k\beta-1}-\frac{\varphi_k\ch\lambda_k\beta}{\ch\lambda_k\beta-1}- \frac{\psi_k\sh\lambda_k\beta}{\lambda_k(\ch\lambda_k\beta-1)}.
\end{gather}
Теперь, подставляя \eqref{o49}--\eqref{o51} в \eqref{o46}, найдем функции $u_k(y)$ в замкнутом виде:
\begin{multline}
\label{o52}
u_k(y)=\varphi_k\frac{\ch\lambda_k\beta-\ch\lambda_ky}{\ch\lambda_k\beta-1}+h_k\frac{\ch\lambda_ky-1}{\ch\lambda_k\beta-1}+ \\ + \psi_k\frac{\sh\lambda_k\beta-\sh\lambda_ky-\sh\lambda_k(\beta-y)}{\lambda_k(\ch\lambda_k\beta-1)}.
\end{multline}




Пусть теперь $\varphi(x)=h(x)=\psi(x)\equiv0$. Тогда из \eqref{o51}, \eqref{o52}, \eqref{o10} и \eqref{o45} при всех $k\in\mathbb{N}$ получим
\begin{equation}
\label{o53}
\int _{0}^{l}u(x,y)\sin\mu_kx\,dx=0,\quad  
\int _{0}^{l}f(x)\sin\mu_kx\,dx=0.
\end{equation}
В силу полноты системы $\{\sin\mu_kx\}_{k\geq1}$ из равенств \eqref{o53} следует, что ${u(x,y)=0}$ почти всюду на $[0,l]$ при любом $y\in[0,\beta]$ и $f(x)=0$ почти всюду на $[0,l]$. 
В~силу условий \eqref{o2} и \eqref{o43} функции $u(x,y)\equiv0$ в $\overline{D}$ и $f(x)\equiv0$ на $[0,l]$.~$\square$

\smallskip



Решение задачи \hyperlink{sab:task4}{4} определяется как сумма рядов
\begin{gather}\label{o54}
u(x,y)=\sum\limits_{k=1}^{\infty}u_k(y)\sin\mu_kx,
\\
\label{o55}
f(x)=\sum\limits_{k=1}^{\infty}f_k\sin\mu_kx.
\end{gather}
Здесь коэффициенты $u_k(y)$ и $f_k$ находятся соответственно по формулам \eqref{o52} и \eqref{o51}.

\smallskip

\hypertarget{sab:lemma3}{}

{\small \sc Лемма 3.} \textit{Для коэффициентов рядов {\em \eqref{o54}} и {\em \eqref{o55}} справедливы оценки}
\begin{gather}
\label{o56}
|u_k(y)|\leq C_{10}\Bigl(|\varphi_k|+|h_k|+\frac{|\psi_k|}{k}\Bigr),
\\
\label{o59}
|u'_k(y)|\leq C_{11}\bigl(k|\varphi_k|+k|h_k|+|\psi_k|\bigr),
\\
\label{o59}
|u''_k(y)|\leq C_{12}\bigl(k^2|\varphi_k|+k^2|h_k|+k|\psi_k|\bigr),
\\
\label{o59}
|f_k|\leq C_{13}\bigl(k^2|\varphi_k|+|h_k|+k|\psi_k|\bigr).
\end{gather}

\smallskip

Справедливость оценок \eqref{o56}--\eqref{o59} непосредственно следует из формул \eqref{o52} и \eqref{o51}.
Исходя из оценок \eqref{o56}--\eqref{o59} легко доказывается следующая

\smallskip
\hypertarget{sab:th6}{}

{\small\sc Теорема 6.} \textit{Если $\varphi(x),$ $ h(x) \in  C^3[0,l],$ $\varphi(0)=\varphi''(0)=h(0)={h''(0)=0,}$ $\psi(x)\in C^2[0,l],$ $\psi(0)=\psi(l)=0,$ то существует единственное решение $u(x,y)$ и $f(x)$ задачи {\em \hyperlink{sab:task4}{4}} и~оно определяется рядами {\em \eqref{o54}} и {\em \eqref{o55},} при этом $u(x,y)\in C^2(\overline{D})$ и $f(x)\in C[0,l]$}.

\smallskip
\hypertarget{sab:th7}{}

{\small\sc Теорема 7.} \textit{Для решения задачи {\em \hyperlink{sab:task4}{4}}  справедливы оценки  устойчивости\/}:
\begin{gather}\label{o60}
\|u(x,y)\|_{L_2[0,l]}\leq C_{14}\left(\|\varphi(x)\|_{L_2[0,l]}+\|h(x)\|_{L_2[0,l]}+\|\psi(x)\|_{L_2[0,l]}\right),
\\
\label{o63}
\|f(x)\|_{L_2[0,l]}\leq C_{15}\left(\|\varphi''(x)\|_{L_2[0,l]}+\|h(x)\|_{L_2[0,l]}+\|\psi'(x)\|_{L_2[0,l]}\right),
\\
\label{o63}
\|u(x,y)\|_{C(\overline{D})}\leq C_{16}\left(\|\varphi'(x)\|_{C[0,l]}+\|h'(x)\|_{C[0,l]}+\|\psi(x)\|_{C[0,l]}\right),
\\
\label{o63}
\|f(x)\|_{C[0,l]}\leq C_{17}\left(\|\varphi'''(x)\|_{C[0,l]}+\|h'(x)\|_{C[0,l]}+\|\psi''(x)\|_{C[0,l]}\right).
\end{gather}


\smallskip

\textit{Д\,о\,к\,а\,з\,а\,т\,е\,л\,ь\,с\,т\,в\,о\/} оценок \eqref{o60}--\eqref{o63} проводится аналогично 
теореме~\hyperlink{sab:th3}{3} с использованием оценок \eqref{o56} и \eqref{o59}.

\smallskip

{\small\sc Замечание 2.} Отметим, что в работе  \cite{sab:21} изучена обратная задача для операторного уравнения второго порядка, из которой можно получить явную формулу решения задачи~\hyperlink{sab:task4}{4}, т.е. ряды \eqref{o54} и \eqref{o55}, которые будут сходиться в~$L_2[0,l]$.

\smallskip


\hypertarget{sab:p2.2}{}

\textbf{2.2.} Пусть $g(y)\not\equiv1$ "--- произвольная непрерывная на $[0,\beta]$ функция. 
В~этом случае аналогично п.~\hyperlink{sab:p2.1}{2.1} можно показать единственность решения задачи~\hyperlink{sab:task4}{4}, но при обосновании существования решения возникают принципиальные трудности в силу роста коэффициентов ряда по закону $e^{\lambda_k\beta}$. 
Поэтому здесь предлагается обратная задача на основании прямой задачи~\hyperlink{sab:task2}{2} или \hyperlink{sab:task3}{3}. Остановимся на задаче~\hyperlink{sab:task2}{2}.

\smallskip

\hypertarget{sab:task5}{}

{\small \sc Задача 5.} Найти функции $u(x,y)$ и $f(x)$, удовлетворяющие условиям (\hyperlink{sab:(2')}{$2'$}), \eqref{o3}--\eqref{o5}, \eqref{o29}, \eqref{o42} и \eqref{o43}. 
Здесь $\varphi(x)$, $h(x)$, $\widetilde{\psi}(x)$ и $g(y)$ "--- заданные достаточно гладкие функции.

\smallskip

Пусть существует решение задачи~\hyperlink{sab:task5}{5}, удовлетворяющее условиям \eqref{o9}. 
Введем функции \eqref{o10}, которые удовлетворяют уравнению \eqref{o13}:
$$
u''_k(y)-\lambda^2_ku_k(y)=F_k(y), 
$$
где
$$
F_k(y)=\frac{2}{l}\int _{0}^{l}F(x,y)\sin\mu_kx\,dx=g(y)f_k.
$$
В этом случае общее решение уравнения \eqref{o13} исходя из формулы \eqref{o15} определяется по формуле
\begin{equation}
\label{o64}
u_k(y)=a_ke^{\lambda_ky}+b_ke^{-\lambda_ky}-\frac{f_k}{2\lambda_k} \left(g_{1k}(y)+g_{2k}(y)\right),
\end{equation}
где $a_k$, $b_k$ и $f_k$ "--- неизвестные постоянные,
\begin{equation}
\label{o65}
g_{1k}(y)=\int _{y}^{\beta}g(t)e^{-\lambda_k(t-y)}\,dt,\quad  
g_{2k}(y)=\int _{0}^{y}g(t)e^{-\lambda_k(y-t)}\,dt.
\end{equation}
Для определения этих коэффициентов воспользуемся условиями \eqref{o16}, \eqref{o32} и \eqref{o47}.
Удовлетворяя функции \eqref{o64} граничным условиям \eqref{o16}, \eqref{o32} и \eqref{o47}, найдем систему
\begin{equation}
\label{o66}
\begin{cases}
      \displaystyle a_k+b_k-\frac{f_k}{2\lambda_k}g_{1k}(0)=\varphi_k, \\
      \displaystyle a_ke^{\lambda_k\beta}+b_ke^{-\lambda_k\beta}-\frac{f_k}{2\lambda_k}g_{2k}(\beta)=h_k,\\
      \displaystyle a_k(1-e^{\lambda_k\beta})+b_k(e^{-\lambda_k\beta}-1)-\frac{f_k}{2\lambda_k} [g_{1k}(0)+g_{2k}(\beta)]=\frac{\widetilde{\psi}_k}{\lambda_k}.
    \end{cases}
\end{equation}
При условии, что при любом $k\in\mathbb{N}$
\begin{equation}\label{o67}
\Delta(k)=g_{1k}(0)+g_{2k}(\beta)=\int _{0}^{\beta}g(t) \bigl[e^{-\lambda_kt}+e^{-\lambda_k(\beta-t)}\bigr]\,dt\neq0
\end{equation}
система \eqref{o66} имеет единственное решение:
\begin{multline}
\label{o68} 
a_k=\frac{h_k-\varphi_k-\widetilde{\psi}_k/\lambda_k}{2(e^{\lambda_k\beta}-1)}+ 
\\+
\frac{(\varphi_k+h_k)(1-e^{-\lambda_k\beta})+(\widetilde{\psi}_k/\lambda_k)(1+e^{-\lambda_k\beta})} {2(e^{\lambda_k\beta}-1)} 
\frac{g_{1k}(0)}{g_{1k}(0)+g_{2k}(\beta)},
\end{multline}  
\begin{equation}\label{o69}
\frac{f_k}{2\lambda_k}=-\frac{(\varphi_k+h_k)(1-e^{-\lambda_k\beta})+(\widetilde{\psi}_k/\lambda_k)(1+e^{-\lambda_k\beta})} {2(g_{1k}(0)+g_{2k})(\beta)},
\end{equation}
\begin{multline}
b_k=\frac{1}{2}\Bigl(\varphi_k+h_k-\frac{\widetilde{\psi}_k}{\lambda_k}\Bigr)+ \frac{e^{\lambda_k\beta}
(\varphi_k-h_k+\widetilde{\psi}_k/\lambda_k )} {2(e^{\lambda_k\beta}-1)}-
\\
\label{o70}
-\frac{e^{\lambda_k\beta}\bigl[(\varphi_k+h_k)(1-e^{-\lambda_k\beta})+ (\widetilde{\psi}_k/\lambda_k)(1+e^{-\lambda_k\beta})\bigr]} {2(e^{\lambda_k\beta}-1)} \frac{g_{1k}(0)}{g_{1k}(0)+g_{2k}(\beta)}.
\end{multline}
Найденные значения коэффициентов по формулам \eqref{o68}--\eqref{o70} подставим в \eqref{o64}. 
Тогда получим
\begin{multline}
u_k(y)=
\varphi_k
\Bigl[\frac{e^{\lambda_k(\beta-y)}-\ch\lambda_ky}{e^{\lambda_k\beta}-1}+ \frac{(1-e^{-\lambda_k\beta})g_{1k}(0)(e^{\lambda_ky}-e^{\lambda_k(\beta-y)})} {2(e^{\lambda_k\beta}-1)(g_{1k}(0)+g_{2k}(\beta))}+
\\
\hspace{5cm}
+\frac{(1-e^{-\lambda_k\beta})(g_{1k}(y)+g_{2k}(y))} {2(g_{1k}(0)+g_{2k}(\beta))}\Bigr]+
\\
+ h_k\Bigl[\frac{\sh\lambda_ky}{e^{\lambda_k\beta}-1}+\frac{(1-e^{-\lambda_k\beta})g_{1k}(0)(e^{\lambda_ky}-e^{\lambda_k(\beta-y)})} {2(e^{\lambda_k\beta}-1)(g_{1k}(0)+g_{2k}(\beta))}+  \hspace{2cm}
\\
\hspace{4.6cm}
+\frac{(1-e^{-\lambda_k\beta})(g_{1k}(y)+g_{2k}(y))} {2(g_{1k}(0)+g_{2k}(\beta))}\Bigr]+ 
\\
+ \frac{\widetilde{\psi}_k}{\lambda_k}\Bigl[-\frac{\sh\lambda_ky}{e^{\lambda_k\beta}-1}+ 
\label{o71}
\frac{(1+e^{-\lambda_k\beta})g_{1k}(0)(e^{\lambda_ky}-e^{\lambda_k(\beta-y)})} {2(e^{\lambda_k\beta}-1)(g_{1k}(0)+g_{2k}(\beta))}+
 \hspace{2cm}
 \\
+
\frac{(1+e^{-\lambda_k\beta})(g_{1k}(y)+g_{2k}(y))} {2(g_{1k}(0)+g_{2k}(\beta))}\Bigr].
\end{multline}

Теперь докажем единственность решения задачи~\hyperlink{sab:task5}{5}. 
Пусть $\varphi(x)\hm=h(x)\hm=\widetilde{\psi}(x)\equiv0$ и выполнены условия \eqref{o67} при всех $k\in\mathbb{N}$. Тогда из \eqref{o69} и \eqref{o71} следует, что при всех $k\in\mathbb{N}$ функции $u_k(y)\equiv0$ и $f_k\equiv0$. Получим равенства \eqref{o53}, из которых уже следует, что $u(x,y)\equiv0$ в $\overline{D}$ и  $f(x)\equiv0$ на $[0,l]$.

Пусть при некотором $k=p\in\mathbb{N}$ нарушается условие \eqref{o67}, т.е. $\Delta(p)=0$. 
Тогда однородная задача~\hyperlink{sab:task5}{5} (где $\varphi(x)=h(x)=\widetilde{\psi}(x)\equiv0$) имеет ненулевое решение:
\begin{gather}\label{o72}
u_p(x,y)=\widetilde{u}_p(y)\sin\mu_px,
\\
\label{o74}
f_p(x)=f_p\sin\mu_px,
\end{gather}
где
\begin{equation}
\label{o74}
\widetilde{u}_p(y)=\frac{f_pg_{1p}(0)(\ch\lambda_p(\beta-y)-\ch\lambda_py)}{2\lambda_p(\ch\lambda_p\beta-1)}- \frac{f_p}{2\lambda_p} \bigl(g_{1p}(y)+g_{2p}(y)\bigr),
\end{equation}
$f_p\neq0$ "--- произвольная постоянная.


Теперь возникает вопрос о существовании нулей выражения $\Delta(p)$. Из \eqref{o67} видно, что если функция $g(t)$ на $[0,\beta]$ не меняет знак, т.е. знакопостоянная, то ${\Delta(k)\neq0}$ при любом $k\in\mathbb{N}$. Если $g(t)$ меняет знак, т.е. имеет нули на $[0,\beta]$, то уравнение $\Delta(p)=0$ может иметь корни. В качестве примера возьмем функцию
$$
g(t)=\sin(at+b),\quad  a, b\in\mathbb{R},\quad  a\neq0.
$$
Тогда выражение $\Delta(p)$ примет вид
$$
\Delta(p)=\frac{2\sqrt{A_p^2+B_p^2}}{\lambda_p^2+a^2}\sin \Bigl(\frac{a\beta}{2}+b\Bigr) \sin\Bigl(\varphi_p-\frac{a\beta}{2}\Bigr).
$$
Здесь
$$
A_p=\lambda_p(1-e^{-\lambda_p\beta}),\quad 
B_p=a(1+e^{-\lambda_p\beta}),\quad  
\varphi_p=\arcsin\frac{A_p}{\sqrt{A_p^2+B_p^2}}.
$$
Отсюда следует, что уравнение $\Delta(p)=0$ имеет счетное множество корней относительно $a\beta\big/2$.


Таким образом, нами установлен следующий критерий единственности решения задачи~\hyperlink{sab:task5}{5}.

\smallskip
\hypertarget{sab:th8}{}



{\small\sc Теорема 8.} \textit{Если существует решение задачи {\em \hyperlink{sab:task5}{5},} удовлетворяющее условиям {\em \eqref{o9},} то оно единственно только тогда$,$ когда выполнены условия {\em \eqref{o67}} при всех $k\in\mathbb{N}.$}

\smallskip

При условии \eqref{o67} решение задачи \hyperlink{sab:task5}{5} формально можно построить в~виде сумм рядов \eqref{o54} и \eqref{o55}, где  коэффициенты уже находятся соответственно по формулам \eqref{o71} и \eqref{o69}. 
Покажем, что суммы этих рядов принадлежат соответственно классам $C^2(\overline{D})$ и $C[0,l]$.

\smallskip

\hypertarget{sab:lemma4}{}

{\small\sc Лемма 4.} \textit{Если функция $g(y)$ непрерывна на $[0,\beta]$ и ${|g(y)|\geq m=\rm const>0},$ то существует постоянная $C_0>0$ такая$,$ что при всех $k\in\mathbb{N}$}
\begin{equation}
\label{o75}
|\Delta(k)|\geq\frac{C_0}{k}.
\end{equation}

\smallskip



\textit{Д\,о\,к\,а\,з\,а\,т\,е\,л\,ь\,с\,т\,в\,о\/}. 
На основании теоремы о среднем выражение $\Delta(k)$ из \eqref{o67} представим в виде
$$
\Delta(k)=\int _{0}^{\beta}g(t)
\bigl[e^{-\lambda_kt}+e^{-\lambda_k(\beta-t)}\bigr]\,dt= 
g(\xi)\frac{2(1-e^{-\lambda_k\beta})}{\lambda_k}, \quad  0<\xi<\beta.
$$
Отсюда следует
$$
|\Delta(k)|\geq\frac{2lm(1-e^{-\lambda_1\beta})}{\pi k\sqrt{1-c(l/(\pi k))^2}}>\frac{C_0}{k},
$$
так как
$$
\sqrt{1-c(l/(\pi k))^2}\leq
\begin{cases}
      \sqrt{1-c(l/\pi)^2} & \text{при}\,\ c\leq0, \\
      \,\,\,\,\,\,1 & \text{при}\,\ 0<c<(\pi/l)^2,
\end{cases}
$$
где постоянная $C_0$ зависит от $l$, $\beta$, $c$ и $m$.~$\square$

\smallskip


\hypertarget{sab:lemma5}{}

{\small\sc Лемма 5.} \textit{Для коэффициентов рядов $u_k(y)$ и $f_k$ справедливы следующие оценки}\/:
\begin{gather}\label{o76}
|u_k(y)|\leq C_{18}\Bigl(|\varphi_k|+|h_k|+\frac{1}{k}|\psi_k|\Bigr),
\\
\label{o79}
|u'_k(y)|\leq C_{19}\left(k|\varphi_k|+k|h_k|+|\psi_k|\right),
\\
\label{o79}
|u''_k(y)|\leq C_{20}\left(k^2|\varphi_k|+k^2|h_k|+k|\psi_k|\right),
\\
\label{o79}
|f_k|\leq C_{21}\left(k^2|\varphi_k|+k^2|h_k|+k|\psi_k|\right).
\end{gather}

\textit{Д\,о\,к\,а\,з\,а\,т\,е\,л\,ь\,с\,т\,в\,о\/}.  
Предварительно оценим следующие выражения:
\begin{gather}
1\geq\frac{e^{\lambda_k(\beta-y)}-\ch\lambda_ky}{e^{\lambda_k\beta}-1}\geq -\frac{\ch\lambda_k\beta-1}{e^{\lambda_k\beta}-1}>-\frac{1}{2},
\\
-1\leq\frac{e^{\lambda_ky}-e^{\lambda_k(\beta-y)}}{e^{\lambda_k\beta}-1}\leq1,
\\
|g_{1k}(y)|\leq\frac{\widetilde{C}_1}{k},\quad  |g_{2k}(y)|\leq\frac{\widetilde{C}_2}{k},
\end{gather}
где $\widetilde{C}_1$ и $\widetilde{C}_2$ "--- положительные постоянные, зависящие от $\beta$, $c$, $l$ и $\max\limits_{0\leq y\leq\beta}|g(y)|$.

Тогда на основании оценки \eqref{o75} из формул \eqref{o71} и \eqref{o69} следует справедливость приведенных оценок \eqref{o76}--\eqref{o79}.~$\square$

\smallskip

Теперь на основании оценок \eqref{o76}--\eqref{o79} нетрудно обосновать справедливость следующих утверждений.

\smallskip
\hypertarget{sab:th9}{}

{\small\sc Теорема 9.} \textit{Пусть функции $\varphi(x)$ и $h(x)$ удовлетворяют условиям теоремы {\em \hyperlink{sab:th6}{6},} $\widetilde{\psi}(x)\in C^2[0,l],$ $\widetilde{\psi}(0)=\widetilde{\psi}(l)=0,$ а функция $g(y)$ удовлетворяет условиям леммы {\em \hyperlink{sab:lemma4}{4}.} 
Тогда существует единственное решение $u(x,y)$ и $f(x)$ задачи~{\em\hyperlink{sab:task5}{5},} которое определяется рядами {\em \eqref{o54}} и {\em \eqref{o55},} при этом коэффициенты этих рядов находятся по формулам {\em \eqref{o71}} и {\em \eqref{o69}} и $u(x,y)\in C^2(\overline{D}),$ $f(x)\in C[0,l].$}

\smallskip
\hypertarget{sab:th10}{}

{\small\sc Теорема 10.} \textit{Для решения задачи {\em \hyperlink{sab:task5}{5}} имеют место оценки  устойчивости {\em \eqref{o60}--\eqref{o63},} но с~другими постоянными$,$ и в этих оценках функцию $\psi(x)$ следует заменить на $\widetilde{\psi}(x).$}

\smallskip

\hypertarget{sab:p2.3}{}

\textbf{2.3.} Пусть теперь $f(x)$ "--- известная, а $g(y)$ "--- неизвестная функции. 
В~этом случае дополнительное условие следует уже задать на отрезке $x=x_0\in(0,l)$, $0\leq y\leq\beta$, а другие краевые условия можно задать, комбинируя граничные условия задач \hyperlink{sab:task1}{1}--\hyperlink{sab:task3}{3}:

\hypertarget{B_1}{}
$$
\begin{matrix}
          u(x,0)=\varphi(x),\quad  0\leq x\leq l,\\[2mm]
          u_y(x,0)-u_y(x,\beta)=\widetilde{\psi}(x),\quad  0\leq x\leq l;
\end{matrix}
\eqno{(B_1)}$$ \hypertarget{B_2}{}
$$
\begin{matrix}
          u_y(x,0)=\psi(x),\quad  0\leq x\leq l,\\[2mm]
          u(x,0)-u(x,\beta)=\widetilde{\varphi}(x),\quad 0\leq x\leq l;
\end{matrix}
\eqno{(B_2)}$$
$$
\begin{matrix}
          u(x,\beta)=h(x),\quad  0\leq x\leq l,\\[2mm]
          u_y(x,0)-u_y(x,\beta)=\widetilde{\psi}(x),\quad 0\leq x\leq l;
\end{matrix}
\eqno{(B_3)}$$
\hypertarget{B_4}{}
$$
\begin{matrix}
          u_y(x,\beta)=\widetilde{h}(x),\quad  0\leq x\leq l,\\[2mm]
          u(x,0)-u(x,\beta)=\widetilde{\varphi}(x),\quad  0\leq x\leq l.
\end{matrix}
\eqno{(B_4)}$$
Здесь остановимся на первом варианте (условие (\hyperlink{B_1}{$B_1$})), т.е. поставим следующую задачу.

\smallskip

\hypertarget{sab:task6}{}

{\small\sc Задача 6.} Найти функции $u(x,y)$ и $g(y)$, удовлетворяющие условиям (\hyperlink{sab:(2')}{$2'$}), \eqref{o3}--\eqref{o5}, \eqref{o29} и, кроме того, условиям
\begin{gather}\label{o80}
u(x_0,y)=d(y),\quad 0\leq y\leq \beta,\quad 0<x_0<l,
\\
\label{o81}
g(y)\in C[0,\beta].
\end{gather}
Здесь $\varphi(x)$, $\widetilde{\psi}(x)$, $f(x)$ и $d(y)$ "--- заданные достаточно гладкие функции, при этом $d(0)=\varphi(x_0)$.

\smallskip

Решение задачи (\hyperlink{sab:(2')}{$2'$}), \eqref{o3}--\eqref{o5}, \eqref{o29}, т.е. задачи~\hyperlink{sab:task2}{2}, нами уже построено в виде суммы ряда \eqref{o31}, где коэффициенты $u_k(y)$ определяются по формулам \eqref{o32}. Поскольку $F(x,y)=f(x)g(y)$, имеем $F_k(y)=f_kg(y)$, следовательно,
$$
F_{1k}(y)=f_kg_{1k}(y),\quad F_{2k}(y)=f_kg_{2k}(y),
$$
где $g_{1k}(y)$ и $g_{2k}(y)$ определены формулами \eqref{o65}. 
С~учетом этих изменений формула \eqref{o32} принимает вид
\begin{multline}
u_k(y)=\varphi_k\frac{\ch\lambda_k(\beta-y)-\ch\lambda_ky}{(\ch\lambda_k\beta-1)}- \widetilde{\psi}_k\frac{\sh\lambda_ky}{\lambda_k(\ch\lambda_k\beta-1)}+
\\
+\frac{f_k}{2\lambda_k}
\Bigl[\frac{g_{1k}(0)(\ch\lambda_k(\beta-y)-\ch\lambda_ky-\sh\lambda_ky)}{(\ch\lambda_k\beta-1)}- \hspace{2cm}
\\
\hspace{2cm}
-
 g_{1k}(y)- 
\frac{g_{2k}(\beta)\sh\lambda_ky}{\ch\lambda_k\beta-1}-g_{2k}(y)\Bigr]=
\\\label{o82}
=\frac{f_k}{2\lambda_k}\bigl[g_{1k}(0)A_k(y)-g_{2k}(\beta)B_k(y)-g_{1k}(y)-g_{2k}(y)\bigr]+C_k(y),
\end{multline}
где
$$
A_k(y)=\frac{\ch\lambda_k(\beta-y)-\ch\lambda_ky-\sh\lambda_ky}{\ch\lambda_k\beta-1},
\quad 
B_k(y)=\frac{\sh\lambda_ky}{\ch\lambda_k\beta-1},
$$
$$
C_k(y)=\varphi_k\frac{\ch\lambda_k(\beta-y)-\ch\lambda_ky}{\ch\lambda_k\beta-1}- \widetilde{\psi}_k\frac{\sh\lambda_ky}{\lambda_k(\ch\lambda_k\beta-1)}.
$$
Теперь удовлетворим ряд \eqref{o31} граничному условию \eqref{o80}:
\begin{equation}\label{o83}
u(x_0,y)=\sum\limits_{k=1}^{\infty}u_k(y)\sin\mu_kx_0=d(y),\,\ 0\leq y\leq \beta,
\end{equation}
здесь $u_k(y)$ находится по формуле \eqref{o82}.

Предварительно первое слагаемое в квадратных скобках в правой части формулы \eqref{o82} преобразуем к виду
\begin{multline}
\widetilde{u}_k(y)=g_{1k}(0)A_k(y)-g_{2k}(\beta)B_k(y)-g_{1k}(y)-g_{2k}(y)=
\\
=A_k(y)\int _{0}^{\beta}g(t)e^{-\lambda_kt}\,dt-B_k(y)\int _{0}^{\beta}g(t)e^{-\lambda_k(\beta-t)}\,dt-
\\
-\int _{y}^{\beta}g(t)e^{-\lambda_k(t-y)}\,dt-\int _{0}^{y}g(t)e^{-\lambda_k(y-t)}\,dt=
\\
=\int _{0}^{y}g(t)\bigl[A_k(y)e^{-\lambda_kt}-B_k(y)e^{-\lambda_k(\beta-t)}-e^{-\lambda_k(y-t)}\bigr]\,dt+
\\
+\int _{y}^{\beta}g(t)\bigl[A_k(y)e^{-\lambda_kt}-B_k(y)e^{-\lambda_k(\beta-t)}-e^{-\lambda_k(t-y)}\bigr]\,dt=
\\
\label{o84}
=\int _{0}^{y}g(t)E_k(t,y)\,dt+\int _{y}^{\beta}g(t)H_k(t,y)\,dt.
\end{multline}
Здесь
$$
E_k(t,y)=A_k(y)e^{-\lambda_kt}-B_k(y)e^{-\lambda_k(\beta-t)}-e^{-\lambda_k(y-t)},
$$
$$
H_k(t,y)=A_k(y)e^{-\lambda_kt}-B_k(y)e^{-\lambda_k(\beta-t)}-e^{-\lambda_k(t-y)}.
$$
Тогда с учетом \eqref{o82} и \eqref{o84} равенство \eqref{o83} принимает вид
\begin{multline}
\sum\limits_{k=1}^{\infty}\frac{f_k}{2\lambda_k}
\biggl(\int _{0}^{y}g(t)E_k(t,y)\,dt+ \int _{y}^{\beta}g(t)H_k(t,y)\,dt\biggr)\sin\mu_kx_0=
\\
\label{o85}
=d(y)- \sum\limits_{k=1}^{\infty}C_k(y)\sin\mu_kx_0=\widetilde{d}(y).
\end{multline}
В силу теоремы~\hyperlink{sab:th3}{3} ряд \eqref{o31} сходится равномерно на замкнутой области $\overline{D}$, поэтому в левой части равенства \eqref{o85}, переставляя операции суммирования и интегрирования между собой, получим интегральное уравнение типа первого рода:
\begin{equation}\label{o86}
\int _{0}^{y}g(t)E(t,y)\,dt+\int _{y}^{\beta}g(t)H(t,y)\,dt=\widetilde{d}(y),
\end{equation}
где
\begin{gather}\label{o87}
E(t,y)=\sum\limits_{k=1}^{\infty}\frac{f_k}{2\lambda_k}E_k(t,y)\sin\mu_kx_0,
\\
\label{o88}
H(t,y)=\sum\limits_{k=1}^{\infty}\frac{f_k}{2\lambda_k}H_k(t,y)\sin\mu_kx_0.
\end{gather}
При этом заметим, что
$$
\widetilde{d}(0)=d(0)-\sum\limits_{k=1}^{\infty}C_{1k}(0)\sin\mu_kx_0=d(0)-\varphi(x_0)=0
$$
и левая часть уравнения \eqref{o86} при $y=0$ равна нулю.

Обе части уравнения \eqref{o86} продифференцируем. В результате получим
\begin{equation}
\label{o89}
g(y)\left[E(y,y)-H(y,y)\right]+\int _{0}^{y}g(t)\frac{\partial E(t,y)}{\partial y}\,dt+ \int _{y}^{\beta}g(t)\frac{\partial H(t,y)}{\partial y}\,dt=\widetilde{d}'(y).
\end{equation}
При этом из соотношений \eqref{o87} и \eqref{o88} следует, что
$$
E(y,y)-H(y,y)=\sum\limits_{k=1}^{\infty}\frac{f_k}{2\lambda_k}\left[E_k(y,y)-H_k(y,y)\right]\sin\mu_kx_0\equiv0,
$$
так как
$$
E_k(y,y)-H_k(y,y)\equiv0.
$$
Тогда, еще раз дифференцируя уравнение \eqref{o89}, будем иметь
\begin{multline}\label{o90}
g(y)\Bigl[\frac{\partial E(t,y)}{\partial y}-\frac{\partial H(t,y)}{\partial y}\Bigr]_{t=y}+
 \int _{0}^{y}g(t)\frac{\partial^2E(t,y)}{\partial y^2}\,dt+ 
\\
+ 
 \int _{y}^{\beta}g(t)\frac{\partial^2H(t,y)}{\partial y^2}\,dt=\widetilde{d}''(y).
\end{multline}
Теперь на основании \eqref{o87} и \eqref{o88} вычислим
\begin{multline}
\Bigl[\frac{\partial E(t,y)}{\partial y}-\frac{\partial H(t,y)}{\partial y}\Bigr]_{t=y}= \sum\limits_{k=1}^{\infty}\frac{f_k}{2\lambda_k}\left[\frac{\partial E_k(t,y)}{\partial y}- \frac{\partial H_k(t,y)}{\partial y}\right]\sin\mu_kx_0=
\\
\label{o91}
=\sum\limits_{k=1}^{\infty}f_k\sin\mu_kx_0=f(x_0).
\end{multline}
Если потребовать что $f(x_0)\neq0$, то из \eqref{o90} и \eqref{o91} получим интегральное уравнение Фредгольма второго рода:
\begin{equation}\label{o92}
g(y)-\lambda\int _{0}^{\beta}g(t)K(t,y)\,dt=\mu(y),
\end{equation}
где
\begin{equation}\label{o920}
K(t,y)=
 \begin{cases}
      \cfrac{\partial^2E(t,y)}{\partial y^2}, & 0\leq t\leq y, \\[2mm]
      \cfrac{\partial^2H(t,y)}{\partial y^2}, & y\leq t\leq \beta,
 \end{cases}
\end{equation}
\begin{equation}\label{o93}
\frac{\partial^2E(t,y)}{\partial y^2}= \sum\limits_{k=1}^{\infty}\frac{f_k}{2\lambda_k}\frac{\partial^2 E_k(t,y)}{\partial y^2} \sin\mu_kx_0=\sum\limits_{k=1}^{\infty}\frac{\lambda_k}{2}f_kE_k(t,y)\sin\mu_kx_0,
\end{equation}
\begin{equation}\label{o94}
\frac{\partial^2H(t,y)}{\partial y^2}= \sum\limits_{k=1}^{\infty}\frac{f_k}{2\lambda_k}\frac{\partial^2 H_k(t,y)}{\partial y^2} \sin\mu_kx_0=\sum\limits_{k=1}^{\infty}\frac{\lambda_k}{2}f_kH_k(t,y)\sin\mu_kx_0,
\end{equation}
$$
\lambda=-1/f(x_0), \quad  \mu(y)=\widetilde{d}''(y)/f(x_0).
$$

Заметим, что функции $E_k(t,y)$ и $H_k(t,y)$ имеют производные любого порядка соответственно на замкнутых областях $0\leq t\leq y\leq \beta$ и $0\leq y\leq t\leq \beta$, и $E_k(y,y)=H_k(y,y),$ а по переменной $k$ ограничены. 
Тогда равномерная сходимость рядов \eqref{o93} и \eqref{o94} зависит от $f_k$, т.е. от гладкости функции $f(x)$ на $[0,l]$. Если функция $f(x)\in C^2[0,l]$ и $f(0)=f(l)=0$, то
$$
f_k=-\frac{1}{\mu^2_k}\frac{2}{l}\int _{0}^{l}f''(x)\sin\mu_kx\,dx=-\frac{f^{(2)}_k}{\mu^2_k}
$$
и ряды \eqref{o93} и \eqref{o94} будут  равномерно сходиться на указанных областях. 
Поэтому ядро $K(t,y)$, определенное формулой \eqref{o920}, непрерывно в замкнутом квадрате $0\leq t,y\leq\beta$. Если $d(y)\in C^2[0,\beta]$ и функции $\varphi(x)$ и $\widetilde{\psi}(x)$ удовлетворяют условиям теоремы~\hyperlink{sab:th3}{3}, то функция $\mu(y)$ непрерывна на $[0,\beta]$. 
Итак, уравнение \eqref{o92} представляет собой интегральное уравнение Фредгольма второго рода с~непрерывным ядром и~непрерывной правой частью. Тогда методом последовательных приближений можно доказать однозначную разрешимость данного уравнения в~классе непрерывных на $[0,\beta]$ функций при $|\lambda|<1/(Ml)$, где $M=\max\limits_{\overline{D}}|K(t,y)|$.
Из теории Фредгольма следует также, что если $\lambda$ не является характеристическим числом ядра $K(t,y)$, то интегральное уравнение \eqref{o92} имеет непрерывное на $[0,\beta]$ решение.

Таким образом, нами доказана следующая

\smallskip
\hypertarget{sab:th11}{}

{\small \sc  Теорема 11.} \textit{Пусть $\varphi(x)\in C^3[0,l],$ $\varphi(0)=\varphi''(0)=\varphi(l)=\varphi''(l)=0,$ 
$\widetilde{\psi}(x)\in C^2[0,l],$ 
$\widetilde{\psi}(0)=\widetilde{\psi}(l)=0,$ 
$d(y)\in C^2[0,\beta],$ 
$d(0)=\varphi(x_0),$ 
$f(x)\in C^2[0,l],$ 
$f(0)=f(l)=0,$ $f(x_0)\neq0.$ 
Тогда единственное решение задачи \em \hyperlink{sab:task6}{6} \em существует  при выполнении одного из следующих условий$:$ 
\begin{itemize}
\item[$1)$] $|f(x_0)|>Ml,$ 
\item[$2)$] число $-f^{-1}(x_0)$ не является характеристическим числом ядра $K(t,y).$
\end{itemize} 
 При этом функция $g(y)$ определяется как решение интегрального уравнения {\em \eqref{o92}} в классе $C[0,\beta],$ после чего функция $u(x,y)$ находится как решение задачи {\em \hyperlink{sab:task2}{2}} по формуле {\em \eqref{o31}} и принадлежит $C^2(\overline{D}).$}

\smallskip

Отметим, что в теореме \hyperlink{sab:th11}{11} условие  $f(x_0)\neq0$ существенно, здесь $\widetilde{x}_0\hm =x_0/l\in(0,1)$. 
Пусть при некоторых $\widetilde{x}_0$ и $n=p\in\mathbb{N}$ выражение ${\sin\pi k\widetilde{x}_0=0}$. 
Тогда для функций $f(x)=\sin\mu_px$ и $g(y)\in C[0,\beta]$ аналогично \eqref{o72}--\eqref{o74} строится ненулевое решение однородной задачи~\hyperlink{sab:task6}{6}  (где $u(0,y)=u(l,y)\equiv0$, $\varphi(x)=\widetilde{\psi}(x)\equiv0$, $d(y)\equiv0$):
\begin{equation}
\label{o95}
u(x,y)=\widetilde{u}_p(y)\sin\mu_px,
\end{equation}
где
$$
\widetilde{u}_p(y)=\frac{f_p}{2\lambda_p}\left[g_{1p}(0)A_p(y)-g_{2p}(\beta)B_p(y)-g_{1p}(y)-g_{2p}(y)\right],
$$
$f_p\neq0$ "--- произвольная постоянная. Заметим, что здесь $f(x_0)=0$.

Следовательно, в силу построенного примера \eqref{o95},   однозначная разрешимость задачи~\hyperlink{sab:task6}{6}   зависит также от выбора точки $\widetilde{x}_0$, т.е. от отношения $x_0/l$. Соотношение $\sin\pi k\widetilde{x}_0=0 \Leftrightarrow \widetilde{x}_0=p/q$, где $p,\,q\in\mathbb{N}$, т.е. тогда, когда отношение $x_0/l$ является рациональным числом из интервала $(0,1)$.

В заключение отметим, что аналогично можно исследовать аналоги задачи~\hyperlink{sab:task6}{6}  на основании граничных условий (\hyperlink{B_2}{$B_2$})--(\hyperlink{B_4}{$B_4$}). В этих случаях относительно неизвестной функции $g(y)$ получится интегральное уравнение типа~\eqref{o92}.

\bigskip

\AdditionalContent{Конкурирующие интересы}{Мы не имеем конкурирующих
интересов.}

\AdditionalContent{Авторский вклад и ответственность}{Все авторы
принимали участие в разработке концепции статьи и в написании
рукописи. Авторы несут полную ответственность за предоставление
окончательной рукописи в печать. Окончательная версия рукописи была
одобрена всеми авторами.}

\AdditionalContent{Финансирование}{Работа выполнена при финансовой
поддержке  Российского фонда фундаментальных исследований (проект №~16--31--00421\_мол\_а),  Российского фонда фундаментальных исследований и~Республики Башкортостан (проект № 17--41--020516\_р\_а).}



%{\thanks{Работа выполнена при поддержке РФФИ-РБ $($проект \rm \No~$)$ и РФФИ $($проект \rm \No~$16-31-00421)$.}}

\begin{thebibliography}{99}

\RBibitem{sab:1} 
\by Сабитов~К.~Б., Мартемьянова~Н.~В.
\paper Нелокальная обратная задача для уравнения смешанного типа
\jour Изв. вузов. Матем.
\yr 2011
\issue 2
\pages 71--85
\mathnet{http://mi.mathnet.ru/ivm7235}
\mathscinet{http://www.ams.org/mathscinet-getitem?mr=2814823}
%\transl
%\jour Russian Math. (Iz. VUZ)
%\yr 2011
%\vol 55
%\issue 2
%\pages 61--74
%\crossref{10.3103/S1066369X11020083}
%\scopus{http://www.scopus.com/record/display.url?origin=inward&eid=2-s2.0-79953009271}

\RBibitem{sab:2} 
\by  Сабитов~К.~Б., Хаджи~И.~А.
\paper Краевая задача для уравнения Лаврентьева--Бицадзе с неизвестной правой частью
\jour Изв. вузов. Матем.
\yr 2011
\issue 5
\pages 44--52
\mathnet{http://mi.mathnet.ru/ivm7302}
\mathscinet{http://www.ams.org/mathscinet-getitem?mr=2920088}
%\transl
%\jour Russian Math. (Iz. VUZ)
%\yr 2011
%\vol 55
%\issue 5
%\pages 35--42
%\crossref{10.3103/S1066369X11050069}
%\scopus{http://www.scopus.com/record/display.url?origin=inward&eid=2-s2.0-80051604675}

\RBibitem{sab:3} 
\by Сабитов~К.~Б., Мартемьянова~Н.~В. 
\paper Обратная задача для уравнения эллиптико-гиперболического типа с~нелокальным граничным условием
\jour Сиб. матем. журн.
\yr 2012
\vol 53
\issue 3
\pages 633--647
\mathnet{http://mi.mathnet.ru/smj2351}
\mathscinet{http://www.ams.org/mathscinet-getitem?mr=2978580}
%\transl
%\jour Siberian Math. J.
%\yr 2012
%\vol 53
%\issue 3
%\pages 507--519
%\crossref{10.1134/S0037446612020310}
%\isi{http://gateway.isiknowledge.com/gateway/Gateway.cgi?GWVersion=2&SrcApp=PARTNER_APP&SrcAuth=LinksAMR&DestLinkType=FullRecord&DestApp=ALL_WOS&KeyUT=000307396200012}
%\scopus{http://www.scopus.com/record/display.url?origin=inward&eid=2-s2.0-84863209961}

\RBibitem{sab:4} 
\by Сабитов~К.~Б., Мартемьянова~Н.~В. 
\paper Обратная задача для уравнения Лаврентьева--Бицадзе, связанная с~поиском элементов правой части
\jour Изв. вузов. Матем.
\yr 2017
\issue 2
\pages 44--57
\mathnet{http://mi.mathnet.ru/ivm9207}
%\transl
%\jour Russian Math. (Iz. VUZ)
%\yr 2017
%\vol 61
%\issue 2
%\pages 36--48
%\crossref{10.3103/S1066369X17020050}
%\isi{http://gateway.isiknowledge.com/gateway/Gateway.cgi?GWVersion=2&SrcApp=PARTNER_APP&SrcAuth=LinksAMR&DestLinkType=FullRecord&DestApp=ALL_WOS&KeyUT=000408829300005}
%\scopus{http://www.scopus.com/record/display.url?origin=inward&eid=2-s2.0-85014896413}

\RBibitem{sab:5} 
\by Сабитов~К.~Б., Мартемьянова~Н.~В. 
\paper Нелокальная обратная задача для уравнения с оператором Лавретьева--Бицадзе в прямоугольной области 
\jour Доклады АМАН
\yr 2013
\vol 15
\issue 2
\pages 73–86
\elib{https://elibrary.ru/item.asp?id=21207843}

\Bibitem{sab:6} 
\by  Sabitov~K.~B., Martemyanova~N.~V.
 \paper Nonlocal boundary value problem for the third order equation of mixed type 
 \jour Contemporary Analysis and Applied Mathematics
 \yr 2015
 \vol 3
 \issue 2
 \crossref{10.18532/caam.52633}
\pages 153–169

\RBibitem{sab:7} 
\by Крайко~А.~Н., Макаров~В.~Е., Пудовиков~Д.~Е.
\paper К построению головной ударной волны при ``обратном'' расчете сверхзвукового течения методом характеристик
\jour Ж.~вычисл. матем. и матем. физ.
\yr 1999
\vol 39
\issue 11
\pages 1889--1894
\mathnet{http://mi.mathnet.ru/zvmmf1589}
\mathscinet{http://www.ams.org/mathscinet-getitem?mr=1728975}
\zmath{https://zbmath.org/?q=an:1014.76038}
\elib{http://elibrary.ru/item.asp?id=13330367}
%\transl
%\jour Comput. Math. Math. Phys.
%\yr 1999
%\vol 39
%\issue 11
%\pages 1814--1819

\RBibitem{sab:8} 
\by Крайко~А.~Н., Пьянков~К.~С. 
\paper Построение профилей и мотогондол, суперкритических в околозвуковом потоке идеального газа
\jour Ж. вычисл. матем. и матем. физ.
\yr 2000
\vol 40
\issue 12
\pages 1890--1904
\mathnet{http://mi.mathnet.ru/zvmmf1412}
\mathscinet{http://www.ams.org/mathscinet-getitem?mr=1830363}
\zmath{https://zbmath.org/?q=an:0998.76042}
%\transl
%\jour Comput. Math. Math. Phys.
%\yr 2000
%\vol 40
%\issue 12
%\pages 1816--1829

\RBibitem{sab:9} 
\by Лавретьев~М.~М., Романов~В.~Г., Шишатский~С.~П.
\book Некорректные задачи математической физики и анализа
\publaddr М.
\publ Наука
\yr 1980
\totalpages 287 

\RBibitem{sab:10} 
\by Лавретьев~М.~М., Резницкая~К.~Г., Яхно~В.~Г.
\book Одномерные обратные задачи математической физики
\publaddr Новосибирск
\publ Наука
\yr 1982
\totalpages 88

\RBibitem{sab:11} 
\by Романов~В.~Г.
\book Обратные задачи математической физики
\publaddr М.
\publ Наука
\yr 1984
\totalpages 264

\RBibitem{sab:12} 
\by Романов~В.~Г., Кабанихин~С.~И.
\book Обратные задачи геоэлектрики
\publaddr М.
\publ Наука
\yr 1991
\totalpages 304 


\RBibitem{sab:13} 
\by  Денисов~А.~М.
\book Введение в теорию обратных задач
\publaddr М.
\publ МГУ
\yr 1994
\totalpages 206

\Bibitem{sab:14} 
\by  Prilepko~A.~I., Orlovsky~D.~G., Vasin~I.~A.
\book Methods for Solving Inverse Problems in Mathematical Physics
\publaddr Boca Raton 
\publ  CRC Press 
\yr 2000
\crossref{10.1201/9781482292985}
\totalpages xiv+709 pp \nofrills
\serial Monographs and Textbooks in Pure and Applied Mathematics
\vol 231

\Bibitem{sab:15} 
\by Isakov~V. 
\book Inverse problems for partial differential equations
\publaddr New York
\publ Springer
\yr 2006
\serial Applied Mathematical Sciences
\vol 127
\crossref{10.1007/0-387-32183-7}
\totalpages xiii+346~pp \nofrills

\RBibitem{sab:16} 
\by Кабанихин~С.~И.
\book Обратные и некорректные задачи
\publaddr Новосибирск
\publ Сибирское научное изд-во
\yr 2009

\RBibitem{sab:17} 
\by Соловьев~В.~В. 
\paper Обратные задачи определения источника для уравнения Пуассона на плоскости
\jour Ж. вычисл. матем. и матем. физ.
\yr 2004
\vol 44
\issue 5
\pages 862--871
\mathnet{http://mi.mathnet.ru/zvmmf845}
\mathscinet{http://www.ams.org/mathscinet-getitem?mr=2096929}
\zmath{https://zbmath.org/?q=an:1130.35134}
%\transl
%\jour Comput. Math. Math. Phys.
%\yr 2004
%\vol 44
%\issue 5
%\pages 815--824

\RBibitem{sab:18} 
\by Соловьев~В.~В. 
\paper Обратные задачи для эллиптических уравнений на плоскости.~I
\jour Дифференц. уравнения
\yr 2006
\vol 42
\issue 8
\pages 1106--1114
\mathnet{http://mi.mathnet.ru/de11545}
\mathscinet{http://www.ams.org/mathscinet-getitem?mr=2294175}
%\transl
%\jour Differ. Equ.
%\yr 2006
%\vol 42
%\issue 8
%\pages 1170--1179
%\crossref{10.1134/S0012266106080118}

\RBibitem{sab:19} 
\by Соловьев~В.~В. 
\paper Обратные задачи для эллиптических уравнений на плоскости.~II
\jour Дифференц. уравнения
\yr 2007
\vol 43
\issue 1
\pages 101–109

\RBibitem{sab:20} 
\by Соловьев~В.~В.
\thesis Обратные задачи для уравнений эллиптического и параболического типов в пространствах Гельдера
\thesisinfo Автореферат диссерт. \dots\ д-ра физ.-мат. наук. 
\publaddr М.
\yr 2014

\RBibitem{sab:21} 
\by Орловский~Д.~Г.
\paper Об одной обратной задаче для дифференциального уравнения второго порядка в~банаховом пространстве
\jour Дифференц. уравнения
\yr 1989
\vol 25
\issue 6
\pages 1000--1009
\mathnet{http://mi.mathnet.ru/de6873}
\mathscinet{http://www.ams.org/mathscinet-getitem?mr=1008642}
\zmath{https://zbmath.org/?q=an:0692.34064}
%\transl
%\jour Differ. Equ.
%\yr 1989
%\vol 25
%\issue 6
%\pages 730--738

\RBibitem{sab:22} 
\by Орловский~Д.~Г. 
\paper К~задаче определения параметра эволюционного уравнения
\jour Дифференц. уравнения
\yr 1990
\vol 26
\issue 9
\pages 1614--1621
\mathnet{http://mi.mathnet.ru/de7277}
\mathscinet{http://www.ams.org/mathscinet-getitem?mr=1080438}
\zmath{https://zbmath.org/?q=an:0712.34016}
%\transl
%\jour Differ. Equ.
%\yr 1990
%\vol 26
%\issue 9
%\pages 1201--1207

\RBibitem{sab:23} 
\by Орловский~Д.~Г.
\paper Обратная задача Дирихле для уравнения эллиптического типа 
\jour Дифференц. уравнения 
\yr 2008
\vol 44
\issue 1
\pages 119–128

\RBibitem{sab:24} 
\by Сабитов~К.~Б. 
\book Уравнения математической физики
\publaddr М.
\publ Физматлит
\yr  2013


\end{thebibliography}

\newpage

\makeentitle

\AdditionalContent{Competing interests}{We have no competing
interests.}
\AdditionalContent{Authors' contributions and responsibilities}{Each
author has participated in the article concept development and in
the manuscript writing. The authors are absolutely responsible for
submitting the final manuscript in print. Each author has approved
the final version of manuscript.}

\AdditionalContent{Funding}{This paper was written with the support of the Russian Foundation for Basic Research (grant no.~16--31--00421\_mol\_a), Russian Foundation for Basic Research and Republic Bashkortostan (grant no.~17--41--020516\_r\_a).}




\begin{thebibliography}{99}

\Bibitem{sab:1l} 
\by Sabitov~K.~B., Martem'yanova~N.~V.
\paper A nonlocal inverse problem for a mixed-type equation
\zmath{1235.35283}
\jour Russian Math. (Iz. VUZ)
\yr 2011
\vol 55
\issue 2
\pages 61--74
\crossref{10.3103/S1066369X11020083}
\scopus{http://www.scopus.com/record/display.url?origin=inward&eid=2-s2.0-79953009271}

\Bibitem{sab:2l} 
\by Sabitov~K.~B., Khadzhi~I.~A.
\paper The boundary-value problem for the Lavrent’ev-Bitsadze equation with unknown right-hand side
\zmath{1236.35096}
\jour Russian Math. (Iz. VUZ)
\yr 2011
\vol 55
\issue 5
\pages 35--42
\crossref{10.3103/S1066369X11050069}
\scopus{http://www.scopus.com/record/display.url?origin=inward&eid=2-s2.0-80051604675}

\Bibitem{sab:3l} 
\by Sabitov~K.~B., Martem'yanova~N.~V.
\paper An inverse problem for an equation of elliptic-hyperbolic type with a nonlocal boundary condition
\zmath{1256.35028}
\jour Siberian Math. J.
\yr 2012
\vol 53
\issue 3
\pages 507--519
\crossref{10.1134/S0037446612020310}
\isi{http://gateway.isiknowledge.com/gateway/Gateway.cgi?GWVersion=2&SrcApp=PARTNER_APP&SrcAuth=LinksAMR&DestLinkType=FullRecord&DestApp=ALL_WOS&KeyUT=000307396200012}
\scopus{http://www.scopus.com/record/display.url?origin=inward&eid=2-s2.0-84863209961}

\Bibitem{sab:4l} 
\by Sabitov~K.~B., Martem'yanova~N.~V.
\paper The inverse problem for the Lavrent'ev–Bitsadze equation connected with the search of elements in the right-hand side
\jour Russian Math. (Iz. VUZ)
\yr 2017
\vol 61
\issue 2
\pages 36--48
\crossref{10.3103/S1066369X17020050}
\isi{http://gateway.isiknowledge.com/gateway/Gateway.cgi?GWVersion=2&SrcApp=PARTNER_APP&SrcAuth=LinksAMR&DestLinkType=FullRecord&DestApp=ALL_WOS&KeyUT=000408829300005}
\scopus{http://www.scopus.com/record/display.url?origin=inward&eid=2-s2.0-85014896413}

\Bibitem{sab:5l} 
\lang In Russian
\by Sabitov~K.~B., Martem'yanova~N.~V.
\paper Nonlocal inverse problem for the equation with the Lavret'ev--Bitsadze operator in a rectangular domain
\jour Doklady AMAN
\yr 2013
\vol 15
\issue 2
\pages 73–86


\Bibitem{sab:6l} 
\by  Sabitov~K.~B., Martem'yanova~N.~V.
 \paper Nonlocal boundary value problem for the third order equation of mixed type 
 \jour Contemporary Analysis and Applied Mathematics
 \yr 2015
 \vol 3
 \issue 2
 \crossref{10.18532/caam.52633}
 \pages 153–169

\Bibitem{sab:7l} 
\by Krajko~A.~N., Makarov~V.~E., Pudovikov~D.~E.
\paper Construction of the bow wave through upstream computation of a supersonic flow by the method of characteristics
\jour Comput. Math. Math. Phys.
\yr 1999
\vol 39
\issue 11
\pages 1814--1819

\Bibitem{sab:8l} 
\by Krajko~A.~N., P'yankov~K.~S.
\paper Construction of airfoils and engine nacelles that are supercritical in a transonic perfect-gas flow
\jour Comput. Math. Math. Phys.
\yr 2000
\vol 40
\issue 12
\pages 1816--1829

\Bibitem{sab:9l} 
\lang In Russian
\by Lavrent'ev~M.~M., Romanov~V.~G., Shishatskii~S.~P.
\book Nekorrektnye zadachi matematicheskoy fiziki i analiza \rm [Ill-posed Problems of Mathematical Physics and Analysis]
\publaddr Moscow
\publ Nauka
\yr 1980
\totalpages 287 

\Bibitem{sab:10l} 
\lang In Russian
\by Lavrent'ev~M.~M., Reznitskaya~K.~G., Yakhno~V.~G.
\book Odnomernye obratnye zadachi matematicheskoy fiziki \rm [One-dimensional Inverse Problems of Mathematical Physics]
\publaddr Novosibirsk
\publ Nauka
\yr 1982
\totalpages 88


\Bibitem{sab:11l} 
\lang In Russian
\by Romanov~V.~G.
\book Obratnye zadachi matematicheskoy fiziki \rm [Inverse Problems of Mathematical Physics]
\publaddr Moscow
\publ Nauka
\yr 1984
\totalpages 264

\Bibitem{sab:12l} 
\lang In Russian
\by Romanov~V.~G., Kabanikhin~S.~I.
\book Obratnye zadachi geoelektriki \rm [Inverse Problems of Geoelectrics]
\publaddr Moscow
\publ Nauka
\yr 1991
\totalpages 304 


\Bibitem{sab:13l} 
\by Denisov~A.~M.
\book  Elements of the Theory of Inverse Problems
\yr 1999
\crossref{10.1515/9783110943252}
\serial Inverse and Ill-Posed Problems Series 
\vol 14
\totalpages viii+272~pp \nofrills
\publaddr Utrecht
\publ VSP Press

\Bibitem{sab:14l} 
\by  Prilepko~A.~I., Orlovsky~D.~G., Vasin~I.~A.
\book Methods for Solving Inverse Problems in Mathematical Physics
\publaddr Boca Raton 
\publ  CRC Press 
\yr 2000
\crossref{10.1201/9781482292985}
\totalpages xiv+709 pp \nofrills
\serial Monographs and Textbooks in Pure and Applied Mathematics
\vol 231

\Bibitem{sab:15l} 
\by Isakov~V. 
\book Inverse problems for partial differential equations
\publaddr New York
\publ Springer
\yr 2006
\serial Applied Mathematical Sciences
\vol 127
\crossref{10.1007/0-387-32183-7}
\totalpages xiii+346~pp \nofrills

\Bibitem{sab:16l} 
\book Inverse and Ill-posed Problems
\yr 2011
\by Kabanikhin~S.~I. 
\crossref{10.1515/9783110224016}
\serial Inverse and Ill-posed Problems. Theory and Applications
\vol 55
\totalpages xvi+459~pp \nofrills
\publaddr Utrecht
\publ VSP Press

\Bibitem{sab:17l} 
\by Solov'ev~V.~V.
\paper Inverse problems of source determination for the Poisson equation on the plane
\jour Comput. Math. Math. Phys.
\yr 2004
\vol 44
\issue 5
\pages 815--824

\Bibitem{sab:18:el} 
\by Solov'ev~V.~V.
\paper Inverse problems for elliptic equations on the plane:~I
\jour Differ. Equ.
\yr 2006
\vol 42
\issue 8
\pages 1170--1179
\crossref{10.1134/S0012266106080118}

\Bibitem{sab:19el} 
\by Solov'ev~V.~V.
\paper Inverse problems for elliptic equations on the plane:~II
\jour Differ. Equ.
\vol 43 
\yr 2007
\issue 1 
\pages 108--117
\crossref{10.1134/s0012266107010119}


\Bibitem{sab:20:en} 
\by Solov'ev~V.~V.
\thesis  Inverse problems for equations of elliptic and parabolic type in H\"older spaces
\thesisinfo  Thesis, Doctor of Phys.-Math. Sci. \nofrills
\publaddr Moscow
\yr 2014

\Bibitem{sab:21:en} 
\by Orlovsky~D.~G.
\paper  On an inverse problem for a second order differential equation in Banach space
\jour Differ. Equ.
\yr 1989
\vol 25
\issue 6
\pages 730--738

\Bibitem{sab:22en} 
\by Orlovsky~D.~G.
\paper Determination of a parameter of an evolution equation
\jour Differ. Equ.
\yr 1990
\vol 26
\issue 9
\pages 1201--1207

\Bibitem{sab:23en} 
\by Orlovsky~D.~G.
\paper Inverse Dirichlet problem for an equation of elliptic type
\zmath{1153.35082}
\jour Differ. Equ.
\vol 44
\issue 1
\pages 124--134 
\yr 2008

\Bibitem{sab:24u} 
\lang In Russian
\by Sabitov~K.~B. 
\book Uravneniia matematicheskoy fiziki  \rm [Equations of mathematical physics]
\publaddr Moscow
\publ Fizmatlit
\yr  2013


\end{thebibliography}



%\end{document}
